\documentclass[11pt,reqno]{amsart}
\usepackage[T1]{fontenc}
\usepackage{lmodern,amsmath,amssymb,mathtools,microtype,booktabs,array,longtable}
\usepackage[margin=1in]{geometry}
\usepackage{graphicx}
\usepackage{xr-hyper}
\usepackage[hidelinks]{hyperref}
\ifdefined\pdfinfoomitdate\pdfinfoomitdate=1\fi
\ifdefined\pdftrailerid\pdftrailerid{}\fi
\ifdefined\pdfsuppressptexinfo\pdfsuppressptexinfo=15\fi
\numberwithin{equation}{section}
\newtheorem{theorem}{Theorem}[section]
\newtheorem{proposition}[theorem]{Proposition}
\newtheorem{lemma}[theorem]{Lemma}
\newtheorem{corollary}[theorem]{Corollary}
\theoremstyle{definition}\newtheorem{definition}[theorem]{Definition}
\newtheorem{example}[theorem]{Example}
\theoremstyle{remark}\newtheorem{remark}[theorem]{Remark}
\newcommand{\E}{\mathbb E}\newcommand{\Prb}{\mathbb P}
\newcommand{\R}{\mathbb R}\newcommand{\one}{\mathbf1}
\newcommand{\Q}{\mathcal Q}\newcommand{\TV}{\mathrm{TV}}
\newcommand{\cp}{\mathrm{cp}}\newcommand{\ev}{\mathrm{ev}}\newcommand{\st}{\mathrm{st}}
\DeclareMathOperator{\tr}{tr}\DeclareMathOperator{\dist}{dist}
\DeclareMathOperator{\law}{Law}\DeclareMathOperator{\rank}{rank}
\DeclareMathOperator{\supp}{supp}
\emergencystretch=2em

\externaldocument[][nocite]{supplement-refs}[General_Theta_Foundations_I_R76_Supplement.pdf]
\hypersetup{pdftitle={General Theta Foundations I: Acquired Geometry and Causal Resource Transfer},pdfauthor={Qian Qi}}
\title[Acquired geometry and causal resource transfer]{General Theta Foundations I:\\Acquired Geometry and Causal Resource Transfer}
\author{Qian Qi}\date{October 2026}
\subjclass[2020]{Primary 60J10, 68Q45; Secondary 20C15, 52A20, 60G40, 62C20, 81P15, 90C22, 94A29}
\keywords{causal experiment, predictive state, positive realization, minimal memory, finite group, quantum instrument, convex duality}
\begin{document}
\begin{abstract}
We study the positive memory required to reproduce complete causal response laws.
For every finite-group target, closed convex constraints on all command
words admit a permutation realization at no greater memory cost. This
reduces optimal uniform approximation over words of unbounded length to
an attained finite optimization and gives robust multiplicity lower bounds.
A sharp simultaneous-positivity inequality determines the complete
feasibility region for bounded centered coefficient matrices with a
common null vector. It yields a three-piece exact mean--contrast memory
law for an unbounded symmetric-group family, together with boundary
rigidity and optimal strategies. The observable group and predictive
dimension stay fixed while the positive resource changes.
For general finite-command quantum instruments, a physical return polytope
and rare continuation branches give one finite simulator for every
adaptive stopped experiment of bounded target expected cost.
The memory depends on continuation-orbit growth; commuting branches give
polylogarithmic bounds. For a fixed badly approximable qubit phase we
prove a matching logarithmic law, with leading constant one, within the
explicit class of physical-state dictionaries with Born terminal decoders.
The inherited stationary initialization theorem separates the counted
pure start from the continuing memory. Complete proofs of the general
causal comparison, reset, optimal discrimination, variational, dual,
spectral and geometric results are retained with their individual hypotheses.
\end{abstract}
\maketitle
\section{Introduction}\label{sec:current-overview}\label{sec:introduction}

A predictive state contains the information in an experimental past that
can affect future laws. Its linear dimension measures one kind of
compression. A controller that stores probabilities on finitely many
labels must also respect positivity, the preparation of its initial
state, and the effect of subsequent observations. These constraints can
make its minimum memory much larger than the predictive dimension.
They also distinguish exact realization, uniform approximation of every
command word, and approximation of experiments stopped at bounded
expected cost. We study these distinctions within the original theory
of causal resource comparison and optimal statistical discrimination.

A terminal response assigns a probability row $b(w)$ to each finite
command word. An instrument response assigns joint probabilities to
command--report histories and includes the state on which a terminal
measurement acts. A classical realization uses one fixed collection
of nonnegative transition and readout tables. Its initial label is pure
and counted, except in explicitly named mixed-start comparisons.
The physical quantum system, an external policy's storage, program bits,
arithmetic precision and sampling work are separate resources. A private
simulator label is never made visible to the external policy.

The central question is how an algebraic or geometric description of
future responses controls their positive implementation. The present
results give complementary answers. A finite-group specification admits
a cost-preserving permutation reduction even when the required laws
have uncertainty or error. A centered pair of nonnegative coefficient
matrices has an explicit necessary and sufficient synthesis boundary,
which gives the complete mean--contrast resource law of an unbounded
nonabelian family. A return polytope and rare continuation branches give
finite approximations of entire stopped experiments, with memory governed
by the growth of the continuation semigroup.

\subsection*{Convex specifications and infinite-word approximation}

Fix a homomorphism $\gamma:A^*\to G$ onto a finite group. For each
$g\in G$ prescribe a closed convex set $K_g\subseteq\Delta(C)$ of
allowed terminal laws. The requirement is imposed on \emph{every} word
with image $g$. Theorem~\ref{thm:r76-convex-normal-form} proves that an
arbitrary stochastic controller on $M$ labels satisfies these constraints
if and only if there exist a finite $G$-set $X$, a preparation $\lambda$
and a stochastic decoder $D$ with
\[
 \lambda P_gD\in K_g\quad(g\in G),\qquad
 |X|+\mathbf1_{\{\lambda\text{ is mixed}\}}\le M.
\]
For a free mixed start the last condition is simply $|X|\le M$.
The original stochastic matrices are not assumed to satisfy group
relations, and the resulting action need not be faithful. The proof
averages identity words, projects onto closed communicating classes,
and shows that the sandwiched action permutes their stationary rows.
Closed convex constraints survive the limiting average. This extends
the retained exact normal form by replacing an exact observable
intertwining with a constraint-preserving projection.

Let $d_X(b)$ be the best uniform total-variation approximation to a
target $b:G\to\Delta(C)$ on the action $X$, and let
$d_X^{\mathrm{pt}}(b)$ restrict the preparation to a point.
Theorem~\ref{thm:r76-minimax-normal-form} gives the optimal error over
all words and all $M$-label pure-start controllers:
\begin{equation}\label{eq:r76-intro-minimax}
 \mathcal E_M^{\mathrm p}(b)
 =\min\left\{\min_{|X|\le M}d_X^{\mathrm{pt}}(b),
              \min_{|X|\le M-1}d_X(b)\right\}.
\end{equation}
Only finitely many actions occur. The minima are attained, depend
Lipschitz-continuously on the target, and are computable by real
quantifier elimination for algebraic data. This is synthesis for a
supplied finite-group target, not an algorithm for unrestricted
isolation or cutpoint questions about a supplied probabilistic automaton.

Theorem~\ref{thm:r76-robust-multiplicity} supplies lower certificates.
Fourier singular values larger than the allowed error force the
corresponding irreducible multiplicities in the core. Below half the
separation between distinct target residuals, any memory saving over
their number must also pay for a mixed initializer. For a prime cyclic
group, fewer than $p$ labels do no better than the best constant law,
whereas $p$ labels are exact (Corollary~\ref{cor:r76-prime-error}).
Approximation thus need not remove the algebraic memory obstruction
when the error must be uniform over words of unbounded length.

\subsection*{A sharp boundary with two positivity constraints}

The group average and the nonconstant response impose distinct
constraints on a probability factorization. For $n\ge5$ put $P=I-J/n$
and consider a probability column $u$ and a real matrix $A$ satisfying
\[
 0\le A\le u\mathbf1^{\mathsf T},\qquad
 \sum_{ij}A_{ij}=np,\qquad PAP=aP,\qquad
 \operatorname{rank}[A\mid u]<n.
\]
The matrices $A$ and $u\mathbf1^{\mathsf T}-A$ share a preparation
constraint. Theorem~\ref{thm:r76-bounded-coefficients} determines exactly
which $(p,a)$ permit such a pair, constructs a common nonnegative factor
of width $n-1$ at every permitted pair of parameters, and classifies the
coefficient pair on the curved boundary. This is an existence theorem;
it does not claim that every given feasible pair has that factor width.

Put
\[
 p_0=\frac1{n-\sqrt n},\qquad p_1=\frac2n,\qquad
 F_n(p)=\begin{cases}
 p(\sqrt n-1)/(\sqrt n+1),&0\le p\le p_0,\\
 ((n-1)p-1)/((n-1)(np-1)),&p_0\le p\le p_1,\\
 (1-p)/(n-1),&p_1\le p\le1.
 \end{cases}
\]
For $a>0$, the answer is $a\le F_n(p)$. The proof combines the retained
singular-matrix extremum with a simultaneous-positivity inequality.
It permits signed null vectors and nonsymmetric coefficient matrices.
The equality analysis forces the extremal coordinate pattern and then
the allowable common factors. On the middle interval both outcome
constraints are active; scaling the first-outcome bound alone would
give an incorrect threshold.

Let $\chi_n$ be the natural permutation character of $S_n$, and set
\[
 \beta_{p,a}(g)=p+a(\chi_n(g)-1),\qquad
 b_{p,a}=(\beta_{p,a},1-\beta_{p,a}).
\]
The exact probability range is
$0\le p-a\le1$ and $0\le p+(n-1)a\le1$.
For $a>0$, Theorem~\ref{thm:r76-mean-contrast} proves
\[
 \kappa_{S_n}(b_{p,a})=
 \begin{cases}n(n-1),&a\le F_n(p),\\n^2,&a>F_n(p),\end{cases}
 \qquad m(b_{p,a})=1+\kappa_{S_n}(b_{p,a}).
\]
Here $\kappa$ is the least positive permutation core and $m$ counts a
pure start. Negative contrast exchanges the outcomes and replaces
$p$ by $1-p$; zero contrast has one-label memory. At every nonzero
contrast the observable group is $S_n$, the residual number is $n!$,
and the predictive dimension is $1+(n-1)^2$. The mean $p$ is an
average response over the group, not a prior on an unknown parameter.
All minimum cores on the low and middle transition boundary are
classified. The middle boundary forces equal orbit weights, while
the low-mean boundary permits specified nonuniform weights.
For $p\ge2/n$, the lower resource suffices throughout the positive-contrast
probability range. Two commands and binary readout realize the family
as a prepared finite-dimensional quantum experiment.

The retained arbitrary-group results explain its representation-theoretic
part: small-positive-contrast memory equals a permutation-module
multiplicity cover (Theorem~\ref{thm:r75-orbit-multiplicity}), and a
single transitive response set is convex exactly at normal stabilizers
(Theorem~\ref{thm:r75-normality-convexity}). These proofs, the exact
R75 return law, the quotient linear programs and their payoff duals
remain current inputs in the supplement.

\subsection*{Complete stopped laws and finite return geometry}

Each observed instrument report changes the conditional state used by
later calls. Theorem~\ref{thm:r76-return-approximation} assumes a finite
physical polytope containing the preparation. Every returning branch
maps the invariant state set into that polytope; continuing branches
have total conditional probability at most $r<1$. Retain all distinct
normalized continuation images of the seed vertices of length less
than $L$, and denote their number by $G(L)$.

One $G(L)$-label mixed-start model, or a counted pure-start model with
at most $G(L)+1$ labels, is exact before the first truncated excursion.
Uniformly over visible adaptive command, stopping and terminal-POVM
policies of target expected calls at most $B$, its complete stopped-law
error is at most
\begin{equation}\label{eq:r76-intro-excursion}
 \min\left\{1,\frac{(1-r)B}{r^{-L}-1}\right\}.
\end{equation}
The bound also holds for the terminal classical--quantum record,
including nontermination mass. The exceptional-run estimate is sharp
as a stopping-probability inequality
(Lemma~\ref{lem:r76-excursion-budget}); this sharpness alone does not
give a lower bound for every approximating controller.

With $b$ commuting continuation maps and $k$ seeds,
$G(L)\le k\binom{L+b-1}{b}$. Other polynomial orbit-growth estimates
give the corresponding power of $\log(B/\delta)$ for accuracy $\delta$.
The hypotheses concern individual normalized branches, and are
verified directly in the application; average mixing is not
substituted for a return condition.

For the retained undamped three-report qubit instrument, a regular
triangle in the real Bloch plane contains every report-$2$ posterior;
report $1$ is a pure reset, and report $0$ is a rotation of probability
$1/8$. For any fixed known alphabet of $a$ commanded angles,
Corollary~\ref{cor:r76-qubit-polylog} gives explicit branch tables with
\[
 M\le1+3\binom{L+a-1}{a},\qquad
 L=\max\left\{1,\left\lceil
       \log_8\left(1+\frac{7B}{8\delta}\right)\right\rceil\right\}.
\]
One angle therefore needs at most $3L+1$ labels, uniformly in its
numerical value. This improves the damping upper for fixed finite
alphabets. The earlier $O(\sqrt{B/\delta})$ construction still applies
simultaneously to the full angle circle. These are different command
quantifiers, not competing assertions about one unrestricted model.

Within the class of fixed physical-state dictionaries with Born terminal
decoders, the one-angle upper has a matching lower. For a fixed badly
approximable rotation phase and fixed $B>0$,
Theorem~\ref{thm:r76-physical-logarithm} proves
\[
 \lim_{\delta\downarrow0}
 \frac{N_{\rm phys}^{\rm m}(\delta;B)}{\log_8(1/\delta)}
 =\lim_{\delta\downarrow0}
 \frac{N_{\rm phys}^{\rm p}(\delta;B)}{\log_8(1/\delta)}=1.
\]
Rare reset--rotation words prepare separated pure states. Joint stopped
record events force distinct physical dictionary points to approximate
those states; zero-call readouts exclude the pure initializer from that
packing. An orbit-prefix return polygon attains the leading constant.
The golden-ratio phase has an explicit arithmetic certificate and five
seed states. The original algebraic rotation has a rational-table upper
with seven seeds, without an asserted badly approximable certificate.
This lower is deliberately stated for physical dictionaries; arbitrary
classical terminal decoders need not admit that representation.

The general stationary-initialization identity remains in the main
article: except for the one-label case, a predictively stationary
response costs exactly one label beyond its least mixed-start
realization (Theorem~\ref{thm:r75-stationary-initializer}). The complete
noisy-instrument laws, exact sector and worst-phase counts, all-angle
dictionary and boundary rigidity remain current in the supplement.
Approximate truncation need not be stationary, so its extra initializer
is an upper construction rather than an automatic lower bound.

\subsection*{Priority, dependencies, and the original causal programme}

Residual equivalence and stochastic automata are classical
\cite{RabinScott59,Rabin63,Heller65,CarlylePaz71}. An exact deterministic
residual theorem is not itself a new positive-realization principle.
Rabin's isolated-cutpoint reduction and approximate matrix-product
questions have different quantifiers from
\eqref{eq:r76-intro-minimax}. Nonnegative stochastic groups
\cite{Schwarz64,Flor69,Brown64,Schwarz62}, finite-group representations
\cite{Serre77}, equivariant factorization and nonnegative-rank geometry
\cite{GPT13,KPT10,BocciCarliniRapallo11,GillisGlineur12} provide
established ingredients. Bounded simplex factorization and shared
complementary-output factors are also prior theory
\cite{VuThanhGillisLecron23}; generic nonnegative-rank thresholds and
nested-polytope transitions have earlier examples \cite{Gillis12}. Real-algebraic minimization of probabilistic
automata is also prior work \cite{KieferWachter14}; the finite-group
projection makes the entire wordwise approximation constraint finite.

Polyhedral simulation and adaptive composition of quantum instruments
have explicit antecedents \cite{MonrasWinter14,MonrasWinter16,Okay26}.
Phase obstructions and noise-dependent classical-memory lower bounds
are treated by Fanizza, Lumbreras and Winter \cite{Fanizza24}.
Waiting-time martingales for patterns are classical \cite{Li80}.
The added claims are the stated convex-specification reduction, the
simultaneous-positivity synthesis boundary, and the complete adaptive
stopped-law construction from return geometry. A targeted literature
search and finite diagnostics do not certify exhaustive priority.

The original actual-history comparison theorem and its common paid
construction remain in the main article. Its witness is evaluated on
the competitor's actual pre-call histories. The paired supplement keeps
the complete arguments on performed-law compactness, variational and
dual attainment, robust tied optimality, stopping rigidity, finite-width
transfer, rank-constrained discrimination, finite-order quantum
measurement, two-stage reset control, memory and communication,
spectral geometry and optimal strategy implementation. The independent
R58 companion retains calibration and nonintervening directional
sensing. Local geometric conclusions keep their local hypotheses;
independent A1, A2 and further Theta objectives remain separate
research problems. The theorem map distinguishes logical prerequisites
of the current article from historical preservation.

\section{Prepared kernels, acquired task laws and observers}\label{sec:data}
All hidden-state, action and report spaces are standard Borel. A fixed unknown index $\theta\in\Theta$ specifies an actual preparation and positive normalized causal kernels. Positivity means nonnegative measures, not an everywhere strictly positive density. Dependence on the entire physical past may be included in the hidden state. A legal policy sees past visible reports and its admitted randomization, but not $\theta$ or a future audit. Iteration of these kernels defines its finite-prefix laws; the consistent countable path law is their kernel extension.

Each call has a cost $c_i\ge c_*>0$. The decision to make the next call is predictable. In the geometric route, the declared productive-call flag is also predictable; in the primitive-instrument route, productive reports are specified outcomes and can be random. Commitment is permitted only after a productive report, and occurs exactly once almost surely in each model. At that time a specified future audit $Z\in H$ is scored, where $H\subset\R^k$ is nonempty, compact and convex. The prediction is committed before that audit is observed. A prospective joint report--audit kernel is defined at every productive call, including rejected calls. The stopping coin is independent of the audit conditional on the entire current visible report and the past. Write $L=\operatorname{diam}(H)^2$; the case $L=0$ is trivial whenever a legal rule exists.

The core cost $C_\tau$ includes initial preparation charged to the core, failed calls and all productive calls through commitment. A separately executed final audit must be added to the raw budget. This convention is fixed before taking a risk infimum. Physical time, stored program constants and temporary computation are not identified with call count.

\subsection*{Three resource interfaces}
A checkpoint procedure may use its whole precommit history, but must output one element of a fixed dictionary of at most $M$ values. This counts its output width only; it is not a claim of $\log_2 M$ total working bits.

An autonomous observer has an alphabet $S$, $|S|\le M$, and one fixed parameter-independent program. At an intercall boundary its only persistent observer memory is its label $s\in S$. Its next command distribution is a fixed function of $s$. A physical port name, elapsed time, subprogram index, random seed or former report is not an extra free input. During a call its current command, current report and fresh coins are available. Whatever persists to another call must be encoded in $S$. A fixed decoder $r:S\to H$ reads only the committed label.

In the \emph{event-stopped} interface a post-report kernel emits $(s',e)\in S\times\{0,1\}$. An event $e=1$ ends the experiment and retains $s'$ for the decoder; no subsequent command is executed. This public commitment event is not a stored bit used in continuing computation. The same label may be used live on one path and terminal on another. In the \emph{state-stopped} interface a fixed map $h_*:S\to\{0,1\}$ determines halt status. The initial label is live, and type-one labels are globally absorbing terminal labels with no executed command rows. Entry into a terminal label is allowed only after a productive report. Thus live and terminal types are disjoint in this interface, and only in this interface by definition.

Write $R_i^\Theta(N,M)$, $i\in\{\cp,\ev,\st\}$, for the infimum of worst-case expected squared loss over programs with robust almost-sure commitment and $\sup_\theta\E_\theta C_\tau\le N$. An empty legal class has value $+\infty$. A program chosen once at random is not a free persistent seed. The ideal measurable class may have fixed real constants; effective implementation is stated separately in Section~\ref{sec:composition}.

\subsection*{A physically acquired comparison law}
Execute a declared prepared comparison experiment and obtain its joint report--audit probability $K_0(dx,dz)$ on $X\times H$. Its report marginal is $\nu$. The law is genuinely preparable and its specification has its own resource account. It need not be a member of $\Theta$. In particular no finite set of calibration data is silently promoted to a uniform certificate. Set
\begin{equation}\label{eq:derived}
 f(x)=\int zK_0(dz\mid x),\quad
 w(x)=\int\|z-f(x)\|^2K_0(dz\mid x),\quad
 d_r(x)=w(x)+\|f(x)-r\|^2.
\end{equation}
Disintegration gives Borel versions with $f\in H$; null sets may be completed arbitrarily. The profile is
\begin{equation}\label{eq:profile}
 \Q_m(b)=\min_{\substack{C\subset H,\ 1\le|C|\le m\\
                  0\le h\le b,\ \int h\,d\nu=1}}
        \int [w+\dist(f,C)^2]h\,d\nu,
 \qquad b\ge1,\ m\ge1.
\end{equation}
We use $+\infty$ for an infeasible cap or width. The quantities $f,w$ are task coordinates of acquired predictive information, not a complete substitute for the predictive quotient. Their derivation and the quotient construction in Appendix~\ref{sec:quotient} precede any use of local dimension.

\subsection*{Lower certificates on actual histories}\label{sec:lower-cert}
Fix \emph{one} actual model $\theta_*\in\Theta$, independent of the competing program. At each actual reached productive pre-call history $h$ and legal command $a$, write $P_{h,a}(dy,dz)$ for its conditional full current report--audit law. Let $P^0_{h,a}$ be a measurable comparison row on the same report space, with a Borel mark $x=\pi_{h,a}(y)$ of law $\nu$ and
\begin{equation}\label{eq:full-moments}
 \E^0[Z\mid y]=f(\pi_{h,a}y),\qquad
 \E^0[\|Z\|^2\mid y]=w(\pi_{h,a}y)+\|f(\pi_{h,a}y)\|^2.
\end{equation}
These are conditional identities given the \emph{full} report, not just its mark. Versions and inequalities are required on actual reachable histories, not only almost everywhere under an unrelated comparison path law.

Fix $0<l_-\le1\le u_-<\infty$. For all Borel $q:Y\to[0,1]$ and all finite-image Borel $r:Y\to H$, suppose predictable nonnegative defects $a_-(h,a),e_-(h,a)$ obey
\begin{align}
 \int q\,d(P_{h,a})_Y
 &\le u_-\int q\,d(P^0_{h,a})_Y+a_-(h,a),\label{eq:lower-mass}\\
 \int q(y)\|z-r(y)\|^2P_{h,a}(dy,dz)
 &\ge l_-\int q(y)\|z-r(y)\|^2P^0_{h,a}(dy,dz)-e_-(h,a).
 \label{eq:lower-task}
\end{align}
Independent current randomization can be adjoined to both rows: integrating inequalities for each fixed random value proves the same inequalities with the same defects. All command-conditioned lower inequalities hold before averaging the chosen command.

Let $A_\tau$ count productive calls. At this same actual model require, for every legal competing procedure,
\begin{equation}\label{eq:effort}
 \E_* C_\tau\ge\kappa_-\E_* A_\tau,\qquad\kappa_->0.
\end{equation}
At budget $N$ suppose, uniformly over all robustly feasible competitors,
\begin{equation}\label{eq:occupation}
 \E_*\sum_{j\le\tau:\ j\ {\rm productive}}a_-(H_{j-1},a_j)\le A_N,
 \qquad
 \E_*\sum_{j\le\tau:\ j\ {\rm productive}}e_-(H_{j-1},a_j)\le E_N.
\end{equation}
These are actual occupation certificates, not properties merely of the proposed upper program. Uniform row defects give $A_N\le(N/\kappa_-)\bar a_-$ and $E_N\le(N/\kappa_-)\bar e_-$. Lemma~\ref{lem:drift} supplies a nonuniform alternative.

\subsection*{One common upper construction}\label{sec:upper-cert}
Independently, let $\mathcal G$ be a parameter-independent program which, from each of its start histories, reaches one productive report. Its conditional expected attempt cost is at most $\kappa_+>0$ on every surviving history in every $\theta\in\Theta$. It uses $B$ live roles and returns to its start label on rejection; no reset of the physical state is assumed. It forms a virtual mark $x$, a current choice $J$ and a permitted export $T_Jr$ of its task readout. For every finite dictionary $C$ assume
\begin{equation}\label{eq:exports}
 |\{T_jr:j\in J,\ r\in C\}|\le D_0|C|,
 \qquad D_0\in\mathbb N.
\end{equation}
Every continuing phase and coin needed for this construction is included in the $B$ live roles. For fixed $0<l_+\le1\le u_+<\infty$ and $a_+,e_+\ge0$, assume for all Borel $q:X\to[0,1]$ and finite-image Borel $r:X\to H$, conditionally on every such history,
\begin{align}
 \E[q(x)]&\ge l_+\int q\,d\nu-a_+,\label{eq:upper-mass}\\
 \E[q(x)\|Z-T_Jr(x)\|^2]
 &\le u_+\int q(x)d_{r(x)}(x)\nu(dx)+e_+.
 \label{eq:upper-task}
\end{align}
The expectation here includes the within-attempt choice $J$. Unlike the lower, the upper is not required to hold for each choice separately. The choices retained after this call are nevertheless charged through the actual export or live role. These hypotheses express one common causal construction, not modelwise optimizers.

\section{Reversible response groups and common controlled memory}\label{sec:r72-reversible}
The fixed-phase results concern one transition matrix. Several commands require one family of matrices on the same retained labels, and checking every command in isolation is insufficient. We give an intrinsic finite-realization criterion for responses whose minimal linear dynamics is reversible and compact, including any finite family of unitary commands on a prepared finite-dimensional quantum system. We then obtain a matching arithmetic memory law for a pair of noncommuting involutions. This is a theorem about exact controlled response realization, not an assertion that a simulator's private label is an observed sufficient statistic for arbitrary feedback optimization.

Positive realization through invariant polyhedral sets is classical \cite{Heller65,MonrasWinter14,ThonJaeger15}. The argument below uses a section of the realizing simplex before passing to the observable quotient; auxiliary hidden-state responses need not belong to the minimal response space. Compact reversibility then turns containment into equality. The resulting finite-group obstruction concerns all command words, rather than a single spectral restriction.

\subsection{The minimal response space}
Let $A$ and $C$ be nonempty finite alphabets of commands and terminal outcomes. A supplied response is a function $b:A^*\to\Delta(C)$. Its residual after a word $w$ is $b^w(z)=b(wz)$, including the empty word $\varnothing$. Put
\[
 V_b=\operatorname{span}_{\R}\{b^w:w\in A^*\},
 \qquad (L_af)(z)=f(az).
\]
Assume $V_b$ is finite dimensional. It is invariant under each $L_a$, since $L_ab^w=b^{wa}$. Call the response \emph{compactly reversible} if the restrictions $L_a|_{V_b}$ are invertible and the closure $G_b$ of their generated group is a compact subgroup of $\operatorname{GL}(V_b)$. This concerns the observable quotient; it does not require the full physical command group to be finite or fully observable.

An exact stochastic realization with $M$ live labels has row-stochastic matrices $Q_a$, one row-stochastic terminal decoder matrix $D$, and the prescribed initial row $e_1^{\mathsf T}$, such that
\begin{equation}\label{eq:r72-response}
 b(a_1\cdots a_n)=e_1^{\mathsf T}Q_{a_1}\cdots Q_{a_n}D.
\end{equation}
Randomized terminal readout is allowed even when all $Q_a$ are deterministic. The target input word is not encoded into an uncounted clock or persistent seed.

\begin{theorem}[Finite realization of compactly reversible responses]\label{thm:r72-reversible}
For a compactly reversible response, the following are equivalent:

\textup{(i)} An exact finite stochastic realization exists.

\textup{(ii)} The residual set $\{b^w:w\in A^*\}$ is finite.

\textup{(iii)} The observable group $G_b$ is finite.

\textup{(iv)} An exact realization with deterministic transitions and a possibly randomized decoder exists.

The least number of deterministic-transition labels is the residual cardinality. This is not in general the least stochastic label count. If an $M$-label stochastic realization exists, then
\[
 |G_b|\le(2^M-1)!.
\]
This last bound is only a finite combinatorial certificate, not a sharp memory law.
\end{theorem}
\begin{proof}
Suppose \eqref{eq:r72-response} holds, and write $r_w=e_1^{\mathsf T}Q_w$. Let $R=\operatorname{span}_{\R}\{r_w\}\subseteq\R^M$. The assignment $\pi(r_w)=b^w$ extends to a well-defined surjective linear map $\pi:R\to V_b$: a relation $\sum_wc_wr_w=0$ implies, after right multiplication by $Q_zD$, the identical relation among $b^w(z)$ for every $z$. Moreover
\[
 \pi(rQ_a)=L_a\pi(r)\qquad(r\in R).
\]
The simplex section $K=R\cap\Delta_M$ is a compact polytope, invariant under right multiplication by every $Q_a$. Its image $P=\pi(K)$ is a compact polytope, contains every residual, and spans $V_b$. It satisfies $L_aP\subseteq P$. It is important that we intersect the simplex with $R$ first: projecting arbitrary hidden-state responses directly into $V_b$ would not be justified.

For an element $T$ of a compact matrix group, positive powers have a subsequence tending to the identity with exponents tending to infinity. One elementary proof takes a convergent subsequence of powers with successively growing index gaps; quotients of two close powers approach the identity. A finite-order element is covered by its period multiples. Thus positive powers also approach $T^{-1}$. Applying this to $L_a$ and using closedness of $P$ gives $L_a^{-1}P\subseteq P$, and hence $L_aP=P$. All elements of $G_b$ preserve $P$. They permute its finitely many vertices, and this permutation action is faithful because those vertices span $V_b$. Consequently $G_b$ is finite. This proves (i)$\Rightarrow$(iii).

For the displayed bound, a vertex of $K$ is determined by its positive support. If two distinct vertices had the same support, their difference could be added and subtracted in sufficiently small multiples while preserving $R$, unit mass and nonnegativity, contradicting extremality. Thus $K$ has at most $2^M-1$ vertices. A linear image of a polytope has no more vertices than the original polytope, so the faithful vertex action proves the bound.

If $G_b$ is finite, the residual orbit is finite. Conversely, if the residual set is finite, each invertible $L_a$ maps it injectively into itself, hence permutes it. Since it spans $V_b$, the generated group and its closure are finite. This proves (ii)$\Leftrightarrow$(iii). Under (ii), use the residuals as labels, send $b^w$ to $b^{wa}$ on input $a$, and use decoder $b^w(\varnothing)$. This realizes $b$ with deterministic transitions. If any such deterministic-transition machine sends $u$ and $v$ to the same label, every suffix yields the same decoder law, so $b^u=b^v$. This proves its minimal count. Finally (iv)$\Rightarrow$(i) is immediate. No assertion that stochastic transitions can always be determinized without increasing size was used.
\end{proof}

\begin{corollary}[Prepared unitary commands]\label{cor:r72-unitary}
Let $U_a$ be finitely many unitary matrices on $\mathbb C^d$, let $\rho$ be a prepared density matrix, and let $\{E_r:r\in C\}$ be a terminal positive-operator-valued measurement. For
\[
 b(w,r)=\tr(E_r\Phi_w(\rho)),\qquad \Phi_a(X)=U_aXU_a^*,
\]
exact finite stochastic realization for all command words is possible if and only if its observable response group is finite. The group is computed after quotienting only those directions invisible to every future commanded terminal test. Its dimension is at most $d^2$.
\end{corollary}
\begin{proof}
Let $W=\operatorname{span}_{\R}\{\Phi_w(\rho)\}$ in the Hermitian matrices. Each conjugation is a Hilbert--Schmidt isometry. Its forward invariant subspace $W$ is also backward invariant, by finite-dimensional invertibility. Let
\[
 N=\{X\in W:\tr(E_r\Phi_zX)=0\text{ for all }z,r\}.
\]
It is invariant under every $\Phi_a$ and hence under its inverse. The orthogonal complement of $N$ in $W$ is invariant as well. Therefore the induced transformations on $W/N$ are orthogonal for the quotient inner product. The map $X\mapsto(z,r\mapsto\tr(E_r\Phi_zX))$ identifies $W/N$ with $V_b$ and intertwines them with $L_a$. Their generated group has compact closure in the orthogonal group. Apply Theorem~\ref{thm:r72-reversible}. A constant measurement can make this induced group trivial even when the physical unitary group is infinite, so finiteness of the full physical group is not a necessary condition.
\end{proof}

\subsection{A matching memory law for two noncommuting commands}
Fix $x\in\R$ with $\sin(2\pi x)\ne0$. On a protected qubit take
\[
 U_0=\begin{pmatrix}1&0\\0&-1\end{pmatrix},\qquad
 U_1=\begin{pmatrix}\cos(\pi x)&\sin(\pi x)\\
                     \sin(\pi x)&-\cos(\pi x)\end{pmatrix},
 \quad \rho=E=|0\rangle\langle0|.
\]
Both commands are involutions. For known $0\le\alpha\le1$, their channels are
\[
 \Phi_{a,\alpha}(X)=\alpha U_aXU_a^*+(1-\alpha)\tr(X)I/2.
\]
Let $M_x^{\rm ctrl}(\alpha)$ be the least stochastic live memory of the whole binary response $w\mapsto\tr(E\Phi_{w,\alpha}\rho)$, using the same matrices and decoder for every word. In the real Bloch plane, with coordinates $(z,x)$, these commands are $\alpha S_0$ and $\alpha S_1$, where $S_0$ reflects the angle $t$ to $-t$ and $S_1$ reflects it to $2\pi x-t$. Their product is $\alpha^2R_{2\pi x}$.

\begin{theorem}[All-word arithmetic law for an involutive command pair]\label{thm:r72-control-envelope}
For $0<\epsilon<1$ and $\alpha=1-\epsilon$, the arithmetic envelope of Section~\ref{sec:r71-arithmetic} gives
\begin{equation}\label{eq:r72-control-envelope}
 \tfrac14\mathfrak A_x(\epsilon)
 \le M_x^{\rm ctrl}(1-\epsilon)
 \le(1+4\pi^2)\mathfrak A_x(\epsilon).
\end{equation}
The upper is one explicit finite stochastic machine valid for every mixed command word, not separately chosen unary machines. The lower ranges over arbitrary common stochastic realizations.

For irrational $x$, the lower and upper logarithmic growth exponents as $\epsilon\downarrow0$ are respectively
\[
 \frac1{\mu(x)+1},\qquad\frac{\mu(x)-1}{2\mu(x)-1},
\]
with values $0,1/2$ at $\mu=\infty$. A uniform cubic-root law holds precisely for badly approximable $x$; otherwise the normalized count has lower limit zero and upper limit infinity.

At $\alpha=1$, the exact count is the reduced denominator $q$ if $x$ is rational, and is infinite otherwise. At $\alpha=0$, it is two. Every strictly damped pair has a finite exact common realization. These statements concern controlled response memory, not a general multi-command boundary-optimal policy theorem.
\end{theorem}
\begin{proof}
Any common realization restricted to the words $(01)^n$ realizes
\[
 \tfrac12[1+\alpha^{2n}\cos(2\pi nx)]
\]
with transition matrix $Q_0Q_1$. Its size is therefore at least
$m_x(1-\alpha^2)\ge\mathfrak A_x(\epsilon(2-\epsilon))/2$
by Theorem~\ref{thm:r71-envelope}. For $1\le c\le2$, the defining maximum of the envelope obeys $\mathfrak A_x(c\epsilon)\ge\mathfrak A_x(\epsilon)/c$ whenever both arguments are in its domain. This follows for each denominator before taking its infimum. Taking $c=2-\epsilon$ proves the lower in \eqref{eq:r72-control-envelope}.

For the upper, the rational multiple-polygon construction in the proof of Theorem~\ref{thm:r71-envelope} gives a regular $l$-gon $P$ on the unit circle, with vertex $v_0=(1,0)$, such that
\[
 l\le(1+4\pi^2)\mathfrak A_x(\epsilon),\qquad
 \alpha R_{2\pi x}v_j\in P\quad\text{for every vertex }v_j.
\]
The first reflection $S_0$ permutes this regular polygon and zero belongs to it. Consequently $\alpha S_0P\subseteq P$, while $S_1=R_{2\pi x}S_0$ gives $\alpha S_1P\subseteq P$. Choose barycentric rows $Q_a(j,k)$ of $\alpha S_av_j$ in the same vertices. Let $d_j=(1+(v_j)_1)/2$. Starting at vertex zero, induction on the actual command word shows that the expected retained vertex equals the protected Bloch vector. Pairing with the first coordinate proves every response identity. The start is one of the counted vertices; there is no extra initialization mixture.

The two constants in \eqref{eq:r72-control-envelope} transfer both exponents and both uniform-law statements from Theorems~\ref{thm:r71-exponents} and \ref{thm:r71-cubic-rigidity}. At zero damping, a rational phase of reduced denominator $q\ge3$ has a common deterministic $q$-gon realization. The word-pair transition has a primitive $q$th-root mode and needs at least $q$ states by the unit-circle spectral obstruction. For irrational $x$ that same obstruction rules out every finite size. At $\alpha=0$ the response is one at the empty word and $1/2$ after every nonempty word, realized by a first label followed by an absorbing fair label; one label cannot change its decoder. This proves all endpoints.
\end{proof}

A fully algebraic example uses $U_1=5^{-1/2}\left(\begin{smallmatrix}2&1\\1&-2\end{smallmatrix}\right)$. Then $e^{2\pi ix}=(3+4i)/5$ is not a root of unity, since its sum with its conjugate is the rational noninteger $6/5$. Each undamped command in isolation has period at most two; their common all-word response has no finite stochastic realization. The arithmetic exponent of this particular angle is not determined here. The example isolates the failure of commandwise testing without inferring a Diophantine class from finitely many coefficients.

If an independent herald with fixed probability $p$ is placed before each command operation, and only failures apply the supplied command, every common realization above preserves the terminal response under every external command schedule and the first-herald cost $1/p$. The commands are actual inputs to that same retained-state transducer. This does not make its private label observable to the caller, allow uncharged midstream quantum measurements, or prove optimality for a controller free to choose an easier command sequence. The unary optimal squared-task interpretation below continues to use its separately proved physical embedding.

\section{Convex specifications and optimal uniform approximation}
\label{sec:r76-convex-specifications}

The exact normal form has a robust extension. The target need not be an
exactly specified probability table, and the stochastic controller need
not satisfy the group relations, even on its observable states. What is
required is a constraint on every word representing a given group element.
The quantifier over all words is essential.

Fix a finite command alphabet $A$, a finite group $G$, and a surjective
homomorphism $\gamma:A^*\to G$. For every $g\in G$ let $K_g$ be a closed
convex subset of the finite probability simplex $\Delta(C)$. Empty sets
are allowed. An $M$-label pure-start controller consists of a prescribed
point row $e_o$, stochastic command matrices $Q_a$, and a row-stochastic
decoder $D_0$. It satisfies the specification if
\begin{equation}\label{eq:r76-specification}
 e_oQ_wD_0\in K_{\gamma(w)}\qquad(w\in A^*).
\end{equation}
For a mixed start, replace $e_o$ by an arbitrary probability row and do
not charge an additional label for preparing that row.

\begin{theorem}[Finite-group convex-specification normal form]
\label{thm:r76-convex-normal-form}
The specification \eqref{eq:r76-specification} has a pure-start
realization on at most $M$ labels if and only if there are a finite
$G$-set $X$, a probability row $\lambda$ on $X$, and a row-stochastic
decoder $D$ such that
\begin{equation}\label{eq:r76-convex-core}
 \lambda P_gD\in K_g\quad(g\in G),\qquad
 |X|+\mathbf1_{\{\lambda\text{ is not a point mass}\}}\le M.
\end{equation}
For mixed starts the corresponding criterion is $|X|\le M$.
The action on $X$ is permitted to have a kernel. In particular no
faithfulness hypothesis is imposed on an approximating action.
\end{theorem}
\begin{proof}
We give the projection argument explicitly because its use of convexity
replaces the exact observable intertwining in
Theorem~\ref{thm:r73-normal-form}.
Let $\mathcal L=\gamma^{-1}(e)$. Select finitely many identity words
$\ell_1,\ldots,\ell_s$ whose positive matrix entries cover every
directed edge occurring in any $Q_\ell$, $\ell\in\mathcal L$, and put
\[
 T=(I+Q_{\ell_1}+\cdots+Q_{\ell_s})/(s+1).
\]
At most $M^2$ words suffice. Each closed communicating class $C_i$ of
$T$ is closed under every identity-word matrix. The positive diagonal
and Lemma~\ref{lem:r73-lazy-convergence} give a stochastic idempotent
$E=\lim_nT^n$. Its image simplex has disjointly supported vertices
$\mu_1,\ldots,\mu_k$, the stationary rows of the closed classes.
Every identity word preserves those classes, and consequently
\begin{equation}\label{eq:r76-sandwich-identity}
 EQ_\ell E=E\qquad(\ell\in\mathcal L).
\end{equation}
For a word $u$ write $F_u=EQ_uE$ on this image simplex. Choose an
inverse word $\bar u$ with $\gamma(\bar u)=\gamma(u)^{-1}$.
Since $T^n$ is a convex combination of identity-word matrices,
\eqref{eq:r76-sandwich-identity} applied to $u\ell\bar u$ and
$\bar u\ell u$ gives
\[
 F_uF_{\bar u}=E=F_{\bar u}F_u.
\]
Thus $F_u$ is an affine automorphism of the image simplex and permutes
its vertices. If $\gamma(u)=\gamma(v)$, multiplication by the common
inverse $F_{\bar u}$ shows $F_u=F_v$. Replacing the middle $E$ by
$T^n$ in $EQ_uEQ_vE$ also gives $F_uF_v=F_{uv}$, since all words
$u\ell v$ have the same group image. This defines a $G$-action $P_g$
on $X=\{1,\ldots,k\}$.

Write $e_oE=\sum_i\lambda_i\mu_i$ and put $D(i,\cdot)=\mu_iD_0$.
For any representative $u$ of $g$,
\[
 \lambda P_gD=e_oEQ_uED_0
             =\lim_{n\to\infty}e_oT^nQ_uT^nD_0.
\]
The last expression before passage to the limit is a convex combination
of $e_oQ_{\ell u\ell'}D_0$ with $\ell,\ell'\in\mathcal L$.
All these probability rows belong to $K_g$. Closedness and convexity
therefore give $\lambda P_gD\in K_g$. This argument includes $u$
empty and makes no claim that the original controller's word laws
factor through $G$.

If the initial label lies in a closed class $C_i$, then $e_oE=\mu_i$
and $\lambda$ is a point mass. Otherwise the $k$ disjoint nonempty
closed classes exclude at least the initial label, so $k+1\le M$.
This proves the pure-start charge. For a mixed row $\eta$, use
$\eta E$ instead and retain $k\le M$.

Conversely, a point preparation on $X$ needs only the permutation
transitions. For a mixed $\lambda$, a further label $\star$ has decoder
$\lambda D$ and, on the first actual command $a$, moves into $X$ with
row $\lambda P_{\gamma(a)}$. Subsequent transitions are permutations.
Its empty and nonempty word laws are respectively $\lambda D$ and
$\lambda P_{\gamma(w)}D$. This gives the required realization and
finishes both start conventions.
\end{proof}

For an exact probability table $b$, let $\kappa_G(b)$ be the least
$|X|$ among factorizations $b(g)=\lambda P_gD$. This is its mixed-start
minimum by the preceding theorem. The notation $m(b)$ retains the
counted pure-start convention; its exact formula and the residual
cardinality $N$ are given in Theorem~\ref{thm:r73-memory-formula}.

For a target $b:G\to\Delta(C)$ write
$\operatorname{TV}(p,q)=\tfrac12\sum_c|p_c-q_c|$ and define
\[
 \mathcal E_M^{\mathrm p}(b)=
 \inf_{\substack{\text{pure-start controllers on at most }M\text{ labels}}}
 \ \sup_{w\in A^*}\operatorname{TV}
           (b(\gamma(w)),e_oQ_wD_0).
\]
Use $\mathcal E_M^{\mathrm m}$ for the mixed-start convention.
For a finite $G$-set $X$ put
\[
 d_X(b)=\min_{\lambda,D}\max_{g\in G}
          \operatorname{TV}(b(g),\lambda P_gD),\qquad
 d_X^{\mathrm{pt}}(b)=\min_{x,D}\max_{g\in G}
          \operatorname{TV}(b(g),e_xP_gD).
\]
The decoder ranges over stochastic rows; compactness makes both minima
attained. Minima over no candidate actions are interpreted as $+\infty$.

\begin{theorem}[Attained finite optimization for infinite-word error]
\label{thm:r76-minimax-normal-form}
For every target $b$, finite $G$, and integer $M\ge1$,
\begin{align}
 \mathcal E_M^{\mathrm m}(b)&=\min_{|X|\le M}d_X(b),\label{eq:r76-mixed-error}\\
 \mathcal E_M^{\mathrm p}(b)&=
 \min\left\{\min_{|X|\le M}d_X^{\mathrm{pt}}(b),
             \min_{|X|\le M-1}d_X(b)\right\}.
 \label{eq:r76-pure-error}
\end{align}
One representative of each finite $G$-set is sufficient, so these are
finite optimizations and their minima are attained. Their values are
$1$-Lipschitz functions of $b$ in $\max_g\operatorname{TV}$.
If $\mathcal E_M^{\mathrm p}(b)>0$, a sufficiently small perturbation of
$b$ still cannot be exactly realized on $M$ pure-start labels; the same
assertion holds for mixed starts. When the target table is algebraic,
the optimum and a realizing finite controller are computable by real
quantifier elimination. This is a decidability assertion, not a
polynomial-time bound or a claim that quantifier elimination was executed.
\end{theorem}
\begin{proof}
Apply Theorem~\ref{thm:r76-convex-normal-form} to the closed convex
balls $K_g=\{p:\operatorname{TV}(p,b(g))\le\delta\}$.
It sends every controller of error at most $\delta$ to an admissible
action without increasing the label charge or the error. Conversely
every displayed action supplies such a controller. The pure constraint
splits precisely into a point preparation with $|X|\le M$ or an
arbitrary preparation with $|X|\le M-1$.
There are finitely many homomorphisms from the fixed finite group into
each symmetric group of degree at most $M$. For each action the allowed
preparations and decoders form compact sets. Hence the finite minima
are attained, including limits in which a mixed preparation becomes a
point mass. The triangle inequality proves the Lipschitz bound and
the perturbation assertion.

For the algorithmic statement, enumerate these actions and introduce
variables $\lambda_x,D_{xc},\delta$ and $u_{gc}\ge0$ with
\[
 -u_{gc}\le b(g,c)-(\lambda P_gD)_c\le u_{gc},\qquad
 \sum_cu_{gc}\le2\delta.
\]
Together with the stochastic constraints these are polynomial
inequalities over the real algebraic numbers. The point-start branch
fixes $\lambda=e_x$ and enumerates $x$. The compact attained minimum
is the least real point of the resulting semialgebraic set of feasible
$\delta$'s. Real quantifier elimination and algebraic sample-point
construction, used with the same scope as
Theorem~\ref{thm:r73-algebraic-synthesis}, compute the value and tables.
\end{proof}

The finite optimization concerns synthesis against a supplied finite-group
target. It does not solve the cutpoint, equivalence, or isolation problem
for a supplied arbitrary probabilistic automaton. Nor does it replace
the specification for all words by a check of one representative per group
element on the original stochastic controller. Those restrictions would
change the problem.

\begin{theorem}[Robust multiplicity and initialization certificates]
\label{thm:r76-robust-multiplicity}
Use the stacked Fourier matrix $A_\rho(b)$ of
\eqref{eq:r75-fourier-block}, with singular values in decreasing order,
and write $m_\rho(X)$ for the multiplicity of $\rho$ in
$\mathbb C[X]$. If an action $X$ approximates $b$ to total-variation
error at most $\delta$, then
\begin{equation}\label{eq:r76-fourier-robust}
 \sigma_{m_\rho(X)+1}(A_\rho(b))\le\sqrt2\,\delta
                       \qquad(\rho\in\widehat G).
\end{equation}
A singular value past the number of columns means zero. For a binary
response $b=(\beta,1-\beta)$ the sharper equivalent scalar certificate is
\begin{equation}\label{eq:r76-binary-robust}
 \sigma_{m_\rho(X)+1}(\widehat\beta(\rho))\le\delta.
\end{equation}
In particular let $r_\rho^\delta$ be the number of singular values of
$A_\rho(b)$ strictly exceeding $\sqrt2\delta$. Then every mixed-start
controller of uniform error at most $\delta$ uses at least
\begin{equation}\label{eq:r76-robust-cover}
 \min\left\{\sum_H q_H[G:H]:q_H\in\mathbb Z_{\ge0},\quad
       \sum_Hq_H\ge1,\quad
       \sum_Hq_Hm_\rho(H)\ge r_\rho^\delta\ (\rho\in\widehat G)
       \right\}
\end{equation}
labels. The scalar version may be used for binary outputs.

Let $N$ be the number of distinct residual tables $b^h(g)=b(hg)$.
If $N>1$, let
\[
 \Delta_b=\min_{b^h\ne b^k}\max_g
                \operatorname{TV}(b(hg),b(kg))>0.
\]
Whenever $\delta<\Delta_b/2$, a pure-start controller with fewer than
$N$ labels must have a projected mixed preparation. Its total budget
is therefore at least one plus the lower bound
\eqref{eq:r76-robust-cover}.
\end{theorem}
\begin{proof}
Put $a(g)=\lambda P_gD$ and $z_c(g)=b(g,c)-a(g,c)$.
Since $\sum_cz_c(g)=0$ and $\tfrac12\sum_c|z_c(g)|\le\delta$,
the sums of the positive and negative parts are each at most $\delta$,
and $\sum_c|z_c(g)|^2\le2\delta^2$. Unitarity gives
\[
 \left\|\begin{pmatrix}z_{c_1}(g)\rho(g^{-1})\\
                  \vdots\\z_{c_{|C|}}(g)\rho(g^{-1})\end{pmatrix}
       \right\|_{\rm op}
 =\left(\sum_c|z_c(g)|^2\right)^{1/2}\le\sqrt2\delta.
\]
Average over $g$ to obtain
$\|A_\rho(b)-A_\rho(a)\|_{\rm op}\le\sqrt2\delta$.
The Fourier-multiplicity lemma gives
$\operatorname{rank}A_\rho(a)\le m_\rho(X)$.
For completeness, if $B$ has rank at most $m$, its nullspace
intersects the span of the first $m+1$ right singular vectors of $A$
in a nonzero vector. On a unit vector in that intersection,
$\|(A-B)v\|=\|Av\|\ge\sigma_{m+1}(A)$.
This proves \eqref{eq:r76-fourier-robust}; the binary version follows
by averaging the scalar error, whose absolute value is at most
$\delta$, against unitary matrices. Summing orbit multiplicities and
using the mixed-start normal form gives the integer bound.

For the last assertion, suppose the projected preparation were a point
$x$. Its distinct residuals are indexed by points $x\cdot h$ in its
orbit. If two target residuals $b^h,b^k$ were represented by the same
point, the two approximating residuals would coincide for every future
$g$, and the triangle inequality would give distance at most $2\delta$
between $b^h$ and $b^k$. This contradicts the definition of $\Delta_b$.
The orbit must therefore contain at least $N$ points. With fewer than
$N$ total labels this is impossible, so the projected start is mixed
and its extra initializer is charged by the normal-form theorem.
\end{proof}

\begin{corollary}[Uniform error cannot replace a prime orbit]
\label{cor:r76-prime-error}
Let $G=C_p$ with $p$ prime and let $b:G\to\Delta(C)$ be any target.
Put $r(b)=\min_{q\in\Delta(C)}\max_g\operatorname{TV}(b(g),q)$.
For either start convention,
\[
 \mathcal E_M(b)=
 \begin{cases}r(b),&1\le M<p,\\0,&M\ge p.\end{cases}
\]
For binary outputs $r(b)=(\max_g\beta(g)-\min_g\beta(g))/2$.
\end{corollary}
\begin{proof}
Every orbit of $C_p$ has size one or $p$. Thus any action with fewer
than $p$ points is trivial and gives a constant response, regardless
of preparation. A one-label controller attains $r(b)$. On $p$ points
the regular action with a point preparation and decoder $b$ is exact.
For binary outputs the midpoint of the smallest interval containing
all values of $\beta$ is a Chebyshev center of that interval.
\end{proof}

For the symmetric-group target \eqref{eq:r75-symmetric-return}, the
scalar standard block is $tI/[n(n-1)]$ and $\Delta_{b_t}=2t/n$.
Indeed a permutation moving $k$ points changes a residual by at most
$kt/n$, with equality at an inverse representative; the smallest
nonzero $k$ is two. The subgroup-index estimate
\eqref{eq:r75-index-estimate} and the preceding theorem imply
\begin{equation}\label{eq:r76-symmetric-robust-lower}
 \mathcal E_M^{\mathrm p}(b_t)\ge\frac{t}{n(n-1)}
       \qquad(n\ge5,\ 0<t\le1,\ M\le n(n-1)).
\end{equation}
If the error were smaller, the full standard multiplicity $n-1$
would be forced, and the small-error initialization charge would give
at least $n(n-1)+1$ labels; here $n!>n(n-1)$.
This explicit lower bound is not asserted to be the optimal error
constant. At every exact resource jump the finite optimization also
gives a strictly positive robustness radius, without a claimed closed
formula for that radius.

% R76 author development: this file extends, and does not replace, the R75 proof.
\section{The complete mean--contrast law and optimal-factor rigidity}
\label{sec:r76-mean-contrast}

The mean of a response and the size of its nonconstant part impose distinct
positivity constraints. The return family of
Theorem~\ref{thm:r75-symmetric-return} fixes the mean at $1/n$.
We now allow every feasible mean and both signs of contrast, for every
$n\ge5$. A second positivity inequality changes the sharp boundary: simply
rescaling the old threshold gives the wrong answer on an interval of means.
The proof yields the complete boundary factorization on that interval.
Throughout, the mean is the average terminal probability over the group;
it is not a prior on an unknown parameter.

Bounded simplex-structured factorization and stacking complementary
nonnegative slacks with a common factor are established techniques
\cite[Section IV-C]{VuThanhGillisLecron23}. Threshold attainment and
boundary identifiability also occur in the geometric theory of nonnegative
rank \cite[Section 4.3]{Gillis12}. Our additional conclusion is the sharp
simultaneous coefficient extremum below, its three-piece parameter region,
and the resulting unrestricted group-memory and boundary-factor laws.

The following elementary extremum treats two nonnegative coefficient
matrices sharing a left null vector. Both outcomes must be represented
through the same preparation; this common-null-vector requirement is
stronger than separate rank bounds on the two outcome matrices.

\begin{lemma}[Sharp simultaneous positivity inequality]
\label{lem:r76-simultaneous-positivity}
Let $n\ge5$ and $v,c,d\in\mathbb R^n$ satisfy
\[
 \sum_i v_i=1,\qquad v\cdot c=v\cdot d=0,\qquad
 c_i\ge v_j-\delta_{ij},\quad d_i\ge\delta_{ij}-v_j
 \quad(1\le i,j\le n).
\]
Put $C=\sum_i c_i$, $D=\sum_i d_i$, and
$\gamma_n=2(n-2)/(n-4)$. For every real $\gamma\ge\gamma_n$,
\begin{equation}\label{eq:r76-weighted-positive-bound}
 \gamma C+D\ \ge\ L_n(\gamma):=
 \frac{\gamma(n+1)+n^2-2n-1+
             2\sqrt{n(\gamma-1)(\gamma+n-2)}}{n-1}.
\end{equation}
Equality holds exactly for the following triples, up to a simultaneous
coordinate permutation. Set
\[
 m=\sqrt{\frac{\gamma+n-2}{n(\gamma-1)}},\qquad
 M=\frac{1+m}{n-1}.
\]
Then
\begin{equation}\label{eq:r76-weighted-equality}
 \begin{split}
 v&=(M,\ldots,M,-m),\\
 c&=(M,\ldots,M,M+M/m),\\
 d&=(1-M,\ldots,1-M,(n-1)M(1-M)/m).
 \end{split}
\end{equation}
In particular $1/\sqrt n<m\le(n-2)/n$. No sign assumption on $v$
or symmetry assumption on the resulting coefficient matrices is made.
\end{lemma}
\begin{proof}
The diagonal inequalities give $D\ge n-1$. If $v\ge0$, the argument
in Lemma~\ref{lem:r75-singular-coefficient} applies verbatim to $c$:
a largest coordinate $M>0$ forces all other positive coordinates to
vanish, hence $C\ge n-1$. We verify below that
$L_n(\gamma)<\gamma(n-1)-1/(n-1)$ in the stated parameter range;
therefore this case is strict.

Otherwise write $m=-\min_i v_i>0$, relabel so that $v_1=M$ is
maximal, let $L$ be the second largest coordinate, and put
$b_0=\max\{L,M-1\}$. The inequalities for $c$ give
\[
 c=(b_0,M,\ldots,M)+u,\quad u\ge0,\qquad
 v\cdot(b_0,M,\ldots,M)=M(b_0+1-M)\ge0.
\]
Using $v\cdot u\ge-m\sum_i u_i$ and $v\cdot c=0$ gives
\begin{equation}\label{eq:r76-C-bound}
 C\ge b_0+(n-1)M+\frac{M(b_0+1-M)}m.
\end{equation}
This is the first-outcome estimate from the retained sharp singular-matrix
proof. For the second outcome, $d=\mathbf1-v+z$ with $z\ge0$ and
$v\cdot z=\|v\|_2^2-1$. Consequently
\[
 D\ge n-1+\frac{1-\|v\|_2^2}{m}.
\]
The bound $-m\le v_i\le M$ implies
$v_i^2\le(M-m)v_i+Mm$; summing it yields
\begin{equation}\label{eq:r76-D-bound}
 D\ge n+\frac{1-M}{m}-nM.
\end{equation}
This bound is valid even when its right side is below $n-1$.

As in the retained proof,
\[
 M\ge M_-:=\frac{1+m}{n-1},\qquad
 L\ge\frac{1-M+m}{n-2},\qquad
 M_+:=1+\frac m{n-1}.
\]
If $M\ge M_+$, \eqref{eq:r76-C-bound} with $b_0\ge M-1$
gives $C\ge nM-1\ge n-1$. Together with $D\ge n-1$ this
case is again strict. If $M_-\le M\le M_+$, substitute
$b_0\ge(1-M+m)/(n-2)$ in \eqref{eq:r76-C-bound}; its right
side increases with $b_0$. This gives
\[
 C\ge F(M):=\frac{1+m}{n-2}+(n-1)M+
                  \frac{n-1}{m(n-2)}(M-M^2).
\]
Therefore
\[
 \gamma C+D\ge G(M):=\gamma F(M)+n+\frac{1-M}{m}-nM.
\]
The quadratic $G$ is strictly concave in $M$, so its minimum on this
interval is at an endpoint. At the left endpoint its value is
\begin{equation}\label{eq:r76-left-endpoint}
 G(M_-)=\frac{\gamma(n+1)+n^2-2n-1+
                  n(\gamma-1)m+(\gamma+n-2)/m}{n-1}
          \ge L_n(\gamma),
\end{equation}
by the arithmetic--geometric mean inequality. At the right endpoint,
\[
 G(M_+)=(\gamma-1)nM_++n-\gamma-\frac1{n-1}
       \ge E_n(\gamma):=\gamma(n-1)-\frac1{n-1}.
\]
For completeness, $E_n(\gamma)>L_n(\gamma)$ throughout the asserted
range. After subtracting the rational terms and multiplying by $n-1$,
this is equivalent to
\[
 n\{(n-3)\gamma-n+2\}
     >2\sqrt{n(\gamma-1)(\gamma+n-2)}.
\]
The left side is positive. On squaring and dividing by $n$, the
difference is
\[
 Q_n(\gamma)=n\{(n-3)\gamma-n+2\}^2
                         -4(\gamma-1)(\gamma+n-2).
\]
For $n\ge6$ its quadratic coefficient is positive,
$Q_n(2)=n(n-6)(n-2)\ge0$, and
$Q_n'(2)=2n(n-3)(n-4)-4n-4>0$. Hence it is strictly positive
at $\gamma\ge\gamma_n>2$. For $n=5$, $\gamma_n=6$ and
$Q_5(\gamma)=16\gamma^2-68\gamma+57$ is positive and increasing
for $\gamma\ge6$. This proves the strict endpoint comparison and
\eqref{eq:r76-weighted-positive-bound}.

Equality forces $M=M_-$ and equality in
\eqref{eq:r76-left-endpoint}; thus $m$ has the stated value. The sum of
the coordinates other than one minimum is $(n-1)M$, so each of them
is $M$. Equality in \eqref{eq:r76-C-bound} places every excess
coordinate of $c$ at the unique negative coordinate and gives
$c_n=M+M/m$. Equality in the bound for $D$ similarly places every
excess over $\mathbf1-v$ at that coordinate, giving the stated $d_n$.
The range $\gamma\ge\gamma_n$ is precisely
$1/\sqrt n<m\le(n-2)/n$, so $M+m\le1$. Thus $d_i=1-M\ge m$
for $i<n$, and $d_n\ge1+m$; all the required inequalities hold.
The displayed $c$ also satisfies its inequalities and $v\cdot c=0$.
Direct substitution verifies $v\cdot d=0$ and equality in the bound.
\end{proof}

Let $\chi_n(g)$ count fixed points in the natural action of $S_n$.
For a mean $p\in[0,1]$ and real contrast $a$, define
\begin{equation}\label{eq:r76-mean-contrast-family}
 \beta_{p,a}(g)=p+a\{\chi_n(g)-1\},\qquad
 b_{p,a}=(\beta_{p,a},1-\beta_{p,a}).
\end{equation}
This is a probability table if and only if
\begin{equation}\label{eq:r76-feasible-mean-contrast}
 0\le p-a\le1,\qquad 0\le p+(n-1)a\le1.
\end{equation}
Both endpoint values occur, at a derangement and at the identity.
Put
\begin{equation}\label{eq:r76-mean-threshold}
 p_0=\frac1{n-\sqrt n},\qquad p_1=\frac2n,\qquad
 F_n(p)=\begin{cases}
 p\,\dfrac{\sqrt n-1}{\sqrt n+1},&0\le p\le p_0,\\[4pt]
 \dfrac{(n-1)p-1}{(n-1)(np-1)},&p_0\le p\le p_1,\\[4pt]
 \dfrac{1-p}{n-1},&p_1\le p\le1.
 \end{cases}
\end{equation}
The three formulas agree at the two meeting points.

\begin{theorem}[Sharp synthesis for bounded centered coefficients]
\label{thm:r76-bounded-coefficients}
Let $n\ge5$, $P=I-J/n$, $0<p<1$ and $a>0$. The following are equivalent:
\begin{enumerate}
\item There are a probability column $u$ and a real matrix $A$ such that
\[
 0\le A\le u\mathbf1^{\mathsf T},\qquad
 \sum_{i,j}A_{ij}=np,\qquad PAP=aP,\qquad
 \operatorname{rank}[A\ \ u]<n.
\]
\item There exists such a pair admitting a common nonnegative factorization
\[
 A=\sum_{i=1}^{n-1}w_i\lambda_i^{\mathsf T}D_i,
 \qquad u=\sum_{i=1}^{n-1}w_i\lambda_i^{\mathsf T},
 \qquad w\in\Delta_{n-1},\quad
 \lambda_i\in\Delta_n,\quad D_i\in[0,1]^n.
\]
\item $a\le F_n(p)$.
\end{enumerate}
Thus joined ordinary rank deficiency and a positive factor of width
$n-1$ have the same parameter-feasibility region in this centered
coefficient problem. This does not assert that every feasible pair
away from the boundary has that nonnegative factorization.

For $0<p\le p_0$ and $a=F_n(p)$, all width-$n-1$ factorizations in (ii), up to factor and simultaneous
coordinate relabeling, are
\begin{equation}\label{eq:r76-low-mean-rigidity}
 \lambda_i=\ell_i,\qquad
 D_i=\frac{np}{(n-1)w_i}\ell_i,\qquad
 \sum_{i=1}^{n-1}w_i=1,\qquad
 w_i\ge\frac{np(1-1/\sqrt n)}{n-1},
\end{equation}
where $\ell_i=(1-1/\sqrt n)e_i+e_n/\sqrt n$ and $D_i$ is the
first-output decoder. For $p_0\le p\le p_1$ and $a=F_n(p)$,
all width-$n-1$ factorizations in (ii) have exactly the form
\begin{equation}\label{eq:r76-middle-rigidity}
 w_i=\frac1{n-1},\qquad
 \lambda_i=(1-b)e_i+b e_n,\qquad
 D_i=e_i+(np-1)e_n,\qquad b=\frac p{np-1},
 \quad 1\le i<n.
\end{equation}
At $p=p_0$ these classifications agree. In the middle interval the
coefficient pair $(A,u)$ itself is unique up to simultaneous coordinate
permutation; the proof gives it explicitly.
\end{theorem}
\begin{proof}
Start with any matrix pair in (i); no positive factorization is assumed
in the converse. Put $B=u\mathbf1^{\mathsf T}-A\ge0$.
A nonzero common left null vector $v$ of $A$ and $u$ also
annihilates $B=u\mathbf1^{\mathsf T}-A$. If its coordinate sum were
zero, $v^{\mathsf T}AP=a v^{\mathsf T}$ would contradict
$v^{\mathsf T}A=0$. Normalize that sum to one.
Exactly as in the retained null-vector decomposition,
\[
 A=a\{I+c\mathbf1^{\mathsf T}-\mathbf1v^{\mathsf T}\},\quad
 B=a\{-I+d\mathbf1^{\mathsf T}+\mathbf1v^{\mathsf T}\},
 \quad v\cdot c=v\cdot d=0.
\]
Nonnegativity gives all the inequalities of
Lemma~\ref{lem:r76-simultaneous-positivity}, and
\begin{equation}\label{eq:r76-mean-null-sums}
 C=\frac pa,\qquad D=\frac{1-p}{a}.
\end{equation}
For $p\le p_0$, apply Lemma~\ref{lem:r75-singular-coefficient} to
$A/(np)$: its centered coefficient is $a/(np)$, so its parameter
there is $a/p$, giving $a\le p\tau_n=F_n(p)$.
For $p\ge p_1$, the diagonal inequalities for $d$ give $D\ge n-1$,
hence $a\le(1-p)/(n-1)=F_n(p)$.

For $p_0<p\le p_1$, set
\[
 m_0=(n-1)p-1,\qquad
 \gamma=\frac{n m_0^2+n-2}{n m_0^2-1}.
\]
Then $1/\sqrt n<m_0\le(n-2)/n$, so
$\gamma\ge\gamma_n$, and $m_0$ is the minimizing value in
\eqref{eq:r76-left-endpoint}. The equality triple
\eqref{eq:r76-weighted-equality} with $M=p$ has sums
\[
 C_0=p(n+1/m_0),\qquad D_0=(1-p)(n+1/m_0).
\]
Thus $L_n(\gamma)=(\gamma p+1-p)(n+1/m_0)$.
Substituting \eqref{eq:r76-mean-null-sums} into the lemma and cancelling
the positive factor $\gamma p+1-p$ gives
\[
 a\le\frac1{n+1/m_0}
     =\frac{(n-1)p-1}{(n-1)(np-1)}=F_n(p).
\]
This proves the necessary boundary everywhere.

We give matching constructions at that boundary. In each case, smaller
$a\ge0$ is obtained on the same core by mixing the displayed decoder
with the constant decoder $p$, with mixing weight $a/F_n(p)$.
For $0<p\le p_0$, use $n-1$ factors with equal weights,
preparation rows $\ell_i$ and decoders $np\ell_i$. The largest decoder
entry is $np(1-1/\sqrt n)\le1$, precisely because $p\le p_0$.
The coefficient matrix is $np A_*$, where $A_*$ is the retained
sharp singular extremizer. Its mean is $p$ and its centered part is
$p\tau_nP=F_n(p)P$, as required.

For $p_0\le p\le p_1$, use \eqref{eq:r76-middle-rigidity}.
Here $0<b<1$ and $0<np-1\le1$. Put $r=np-1$ and
$a_*=(1-b)/(n-1)=F_n(p)$. The resulting coefficient matrix has
\[
 A_{ii}=a_*\ (i<n),\quad A_{ij}=0\ (i\ne j<n),\quad
 A_{in}=a_*r,\quad A_{ni}=\frac b{n-1},\quad A_{nn}=br=p.
\]
The identity $r(nb-1)=1$ shows that all cross differences of
$A-a_*I$ vanish. Thus $PA P=a_*P$ and the group mean is $p$;
the response is $\beta_{p,a_*}$.

For $p_1\le p<1$, put $b=2/n$, use the same preparations
$(1-b)e_i+b e_n$ and equal orbit weights, and let
\[
 d_0=\frac{np-2}{n-2},\qquad
 D_i(j)=d_0+(1-d_0)\mathbf1_{\{j\in\{i,n\}\}}.
\]
These are valid decoders. At $p=p_1$ they agree with the preceding
construction; mixing those decoders with the constant one shows that
the response is $p+[(1-p)/(n-1)](\chi_n-1)$ for every $p\ge p_1$.
Equivalently this follows by direct permutation summation. This proves
attainment at $F_n(p)$. The same
constructions and decoder mixing give (iii)$\Rightarrow$(ii), while
the displayed factorization immediately gives (ii)$\Rightarrow$(i).
The converse already proved (i)$\Rightarrow$(iii), so the matrix
synthesis equivalence is complete.

For the first rigidity statement, the equality case of the retained
singular-coefficient lemma forces $A=np A_*$. The positive diagonal
leading block and its zero off-diagonal entries force precisely one
nonnegative rank-one summand at each of its $n-1$ diagonal positions.
The adjacent last-row and last-column ratios then give
\eqref{eq:r76-low-mean-rigidity}, exactly as in the retained proof.
Its decoder bound is $w_i\ge np(1-1/\sqrt n)/(n-1)$.
At $p=p_0$ summing these lower bounds forces equal weights. Pair
uniqueness at this endpoint also follows without presupposing a positive
factorization. The singular lemma fixes $A=np_0 A_*$, whose normalized
left null vector is $(p_0,\ldots,p_0,-1/\sqrt n)$. Thus
$v\cdot u=0$ and $\sum_i u_i=1$ force $u_n=1/\sqrt n$.
The bounds $u_i\ge A_{ii}=1/(n+\sqrt n)$ for $i<n$ then sum to
$1-1/\sqrt n$, forcing equality in every one of them.

If $p_0<p\le p_1$ and $a=F_n(p)$, equality holds in the weighted
lemma with the chosen $\gamma$. It forces
$v=(p,\ldots,p,-m_0)$ and the unique $c,d$ in
\eqref{eq:r76-weighted-equality}. The coefficient matrix is therefore
exactly the one displayed in the middle construction, while
$u=a(c+d)$ has $u_i=a$ for $i<n$ and $u_n=b$. Its leading
block again forces one rank-one summand at each diagonal position, with
both factors supported on $\{i,n\}$. The last-row ratios force
$\lambda_i=(1-b)e_i+b e_n$. The diagonal entries force
$D_i(i)=a/\{w_i(1-b)\}=1/\{(n-1)w_i\}$.
The upper bound $D_i(i)\le1$ implies $w_i\ge1/(n-1)$;
summing the weights makes every one equal to $1/(n-1)$.
The last-column ratios then give $D_i(n)=np-1$.
This proves the complete middle-boundary classification.
\end{proof}

\begin{corollary}[Exact memory throughout the mean--contrast plane]
\label{thm:r76-mean-contrast}
Let $n\ge5$ and let $(p,a)$ satisfy
\eqref{eq:r76-feasible-mean-contrast}. At $a=0$, both mixed-core and
pure-start minima are one. If $a\ne0$, put $q=p$ when $a>0$ and
$q=1-p$ when $a<0$. Then
\begin{equation}\label{eq:r76-full-memory-law}
 \kappa_{S_n}(b_{p,a})=
 \begin{cases}
 n(n-1),&0<|a|\le F_n(q),\\
 n^2,&|a|>F_n(q),
 \end{cases}
 \qquad m(b_{p,a})=1+\kappa_{S_n}(b_{p,a}).
\end{equation}
Every nonconstant table has observable group $S_n$, residual cardinality
$n!$, and observable dimension $1+(n-1)^2$.
For $a>0$ and $p\ge2/n$, the smaller memory realizes every feasible
contrast. Thus positivity causes no additional core cost anywhere in
that part of the parameter plane.

At the nontrivial boundary $a=F_n(p)$ with $0<p\le2/n$, all
minimum permutation-core factorizations are exactly those in
\eqref{eq:r76-low-mean-rigidity} and \eqref{eq:r76-middle-rigidity}.
For negative contrast, exchange the two output labels. This classifies
minimum permutation cores, without asserting uniqueness of every
unreduced stochastic machine.
\end{corollary}
\begin{proof}
It suffices to treat $a>0$, because exchanging the two outputs replaces
$(p,a)$ by $(1-p,-a)$. The trivial block of the Fourier transform is
$(p,1-p)$, and the nontrivial standard block is
$aI_{n-1}/(n-1)$ in the first coordinate. The faithful-standard-module
argument in Theorem~\ref{thm:r75-symmetric-return} therefore gives the
claimed observable group, residual count and dimension for every $a>0$.
That theorem's subgroup-index estimate depends only on this module.
It implies the lower bound $n(n-1)$ and proves that every core of size
less than $n^2$ compresses to exactly $n-1$ natural orbits: extra orbits
have no standard component, so their nonnegative constant contribution
is absorbed into the natural-orbit decoders. This absorption preserves
both output probabilities and does not use the value $p=1/n$.

On $n$ natural orbits, put weight $1/n$ on each, prepare point $i$
on orbit $i$, and use
\[
 D_i(j)=p-a+na\,\mathbf1_{\{i=j\}}.
\]
These decoders lie in $[0,1]$ exactly under
\eqref{eq:r76-feasible-mean-contrast}. Their response is
\eqref{eq:r76-mean-contrast-family}. Thus the core is always at most
$n^2$, and there are no possible intermediate values.

Suppose a realization on $n-1$ natural orbits exists. Its preparation
rows $\lambda_i$, weights $w_i$ and decoder rows $D_i$ give
\[
 A=\sum_{i<n}w_i\lambda_i^{\mathsf T}D_i,\qquad
 u=\sum_{i<n}w_i\lambda_i^{\mathsf T},\qquad
 B=u\mathbf1^{\mathsf T}-A\ge0.
\]
Here $A\ge0$, $u$ is a probability column, and both coefficient
matrices have the same preparation factor. Their permutation sums give
the two outcomes. The permutation-switch argument in the retained proof
implies $PAP=aP$, and averaging gives $\sum_{ij}A_{ij}=np$.
The preparation matrix has at most $n-1$ rows, so
$\operatorname{rank}[A\ \ u]<n$.
The matrix criterion in
Theorem~\ref{thm:r76-bounded-coefficients} is therefore necessary.
Conversely its width-$n-1$ positive factorizations give exactly
$n-1$ natural-orbit preparations and valid binary decoders.
Their coefficient identities imply the required response for every
permutation. Hence a core below $n^2$ exists exactly when $a\le F_n(p)$,
and then the module lower bound makes its size $n(n-1)$.
The retained pure-start formula gives $m=\min\{n!,1+\kappa\}$,
which equals $1+\kappa$ because $n^2+1<n!$ for $n\ge5$.
At $a=0$ the constant decoder uses one state. Every minimum core on the
nontrivial boundary consists of $n-1$ natural orbits, so the factor
classification follows from the matrix theorem.
\end{proof}

\paragraph{Physical realization and scope.}
Use the prepared unitary experiment of
Theorem~\ref{thm:r75-symmetric-return}, with
$\rho=n^{-1}\sum_i|ii\rangle\langle ii|$ and
$E=\sum_i|ii\rangle\langle ii|$. Replacing the terminal effect by
\[
 E_{p,a}=(p-a)I+naE
\]
produces \eqref{eq:r76-mean-contrast-family}; its two eigenvalues are
$p-a$ and $p+(n-1)a$, so the exact validity range is
\eqref{eq:r76-feasible-mean-contrast}, including negative contrast.
The same two command generators realize every word.
The theorem solves the whole displayed two-parameter family, including
its constant boundary and its minimum-core factorization on the
nontrivial transition curve. It is not an unrestricted high-contrast
formula for every finite group, an identification result for an unknown
model, or a solution of the original arbitrary-prior decision-risk problem.

\section{The exact cost of initializing a stationary response}
\label{sec:r75-stationary-initialization}

The distinction between a freely chosen initial distribution and a
prescribed initial label is not peculiar to a periodic example.  There is
an exact general law.  Its proof is an elementary elimination in the
space of future response functions, within the classical positive
realization framework \cite{Heller65,MonrasWinter14,MonrasWinter16}.  It does not require
hidden states to lie in the minimal linear realization, or require their
unobserved transition matrix to be irreducible.

Let $A$ be a nonempty finite command alphabet, with a finite report set
$Y_a$ for each $a\in A$.  A finite classical realization consists of
nonnegative matrices $T_{a,y}$ such that
\[
 T_a:=\sum_{y\in Y_a}T_{a,y},\qquad T_a\one=\one,
\]
and an initial probability row $\mu$.  For a command--report word
$w=(a_1,y_1)\cdots(a_n,y_n)$ write
$T_w=T_{a_1,y_1}\cdots T_{a_n,y_n}$ and
$p(w)=\mu T_w\one$.  Optional terminal readouts have decoder
columns $d_{b,z}\ge0$, with $\sum_zd_{b,z}=\one$, and response
$p(w;b,z)=\mu T_wd_{b,z}$.  All following equalities include every specified
terminal readout when these are part of the experiment.

Call the prepared response \emph{predictively stationary under $\nu$}
if $\nu$ is a probability row on $A$ with every $\nu(a)>0$ and
\begin{equation}\label{eq:r75-predictive-stationarity}
 \sum_{a\in A}\nu(a)\sum_{y\in Y_a}p((a,y)w)=p(w)
 \qquad\text{for every future word }w.
\end{equation}
When terminal readouts are specified, additionally require
\begin{equation}\label{eq:r75-terminal-stationarity}
 \sum_{a\in A}\nu(a)\sum_{y\in Y_a}p((a,y)w;b,z)=p(w;b,z)
 \qquad\text{for all }w,b,z,
\end{equation}
including the empty word $w$.
This is an equality of observable responses; it does not assume
$\mu\sum_a\nu(a)T_a=\mu$ in a particular hidden realization.  For one
command and no terminal readout it is ordinary stationarity.  A response
has a one-label realization precisely when its report probabilities
factor as $p(w)=\prod_jp_{a_j}(y_j)$, and every terminal readout has a
fixed distribution independent of the preceding word.  For one command
this is the i.i.d. case.

If a stationary hidden row is desired, it can be chosen without adding
labels. Put $\overline T=\sum_a\nu(a)T_a$. Every $\mu\overline T^j$
has the same future response by the displayed stationarity identities.
A subsequential limit of the Ces\`aro averages
$H^{-1}\sum_{j=0}^{H-1}\mu\overline T^j$ is stationary for
$\overline T$ and retains that response. Thus the freely mixed minimum
below agrees with the convention requiring a stationary hidden row for
the averaged command, and with the usual stationary HMM convention when
there is only one command.

Let $n_{\mathrm{mix}}(p)$ be the minimum number of labels with an arbitrary
initial probability row, and let $n_{\mathrm{pure}}(p)$ be the minimum when
one prescribed initial label is counted.  The minima are $+\infty$ when
there is no finite realization.

\begin{theorem}[Stationary initialization law]
\label{thm:r75-stationary-initializer}
Suppose \eqref{eq:r75-predictive-stationarity} holds, together with
\eqref{eq:r75-terminal-stationarity} when terminal readouts are included.
If the response has
a one-label realization, both minima are one.  Otherwise
\begin{equation}\label{eq:r75-initialization-law}
 n_{\mathrm{pure}}(p)=n_{\mathrm{mix}}(p)+1,
 \qquad (+\infty)+1:=+\infty.
\end{equation}
The lower holds for all real nonnegative transition tables, including
reducible and nonminimal linear realizations.  The upper processes the
first actual command while entering the core; it neither discards a
physical call nor inserts an uncounted preliminary randomization.
\end{theorem}
\begin{proof}
Take any finite pure-start realization, with initial label $o$, and let
$f_i$ denote the complete future response from hidden label $i$.  Thus
$f_i(w)=e_i^{\mathsf T}T_w\one$, with the analogous values for terminal
readouts.  Put $\overline T=\sum_a\nu(a)T_a$.  Stationarity gives the
functional identity
\begin{equation}\label{eq:r75-initial-functional-identity}
 f_o=\sum_j\overline T_{oj}f_j.
\end{equation}
If $\overline T_{oo}=1$, nonnegativity and full support of $\nu$ imply
$(T_a)_{oo}=1$ for every $a$.  All report transitions from $o$ then return
to $o$.  Their probabilities, and the terminal readouts at $o$, constitute
a one-label realization.  In the case under consideration this is
excluded.  Hence
\[
 \beta_j=\frac{\overline T_{oj}}{1-\overline T_{oo}}\quad(j\ne o)
\]
is a probability row and $f_o=\sum_{j\ne o}\beta_jf_j$.

Delete label $o$.  On the remaining labels define
\begin{equation}\label{eq:r75-elimination-tables}
 \widehat T_{a,y}(i,j)
   =T_{a,y}(i,j)+T_{a,y}(i,o)\beta_j,
 \qquad i,j\ne o,
\end{equation}
and retain their terminal decoders.  These matrices are nonnegative and
their sums over $y$ are stochastic.  The recursion
$f_i((a,y)w)=\sum_jT_{a,y}(i,j)f_j(w)$, with the functional identity for
$f_o$ substituted on its right side, is exactly the recursion for
\eqref{eq:r75-elimination-tables}.  Induction on word length, starting with
the empty word and each terminal decoder, proves that every retained
$f_i$ is unchanged.  Initial row $\beta$ therefore realizes $f_o$ with
one fewer label.  This proves $n_{\mathrm{mix}}\le
n_{\mathrm{pure}}-1$ whenever a finite pure realization exists.

Conversely, given a $k$-label mixed-start realization $(\mu,T_{a,y})$,
append one label $o$.  Its report transition on command $a$ enters the
core with row $\mu T_{a,y}$ and never returns to $o$; the core retains its
original transitions.  A terminal readout at $o$ has row $\mu d_{b,z}$.
The first physical command is already included in this row, so induction
gives the original law at every word length, including zero.  This uses
$k+1$ labels and proves the reverse inequality.  A finite pure realization
is itself a finite mixed one, so the two finiteness assertions are
equivalent.  The one-label exception is immediate.
\end{proof}

\begin{corollary}[Stationary processes]
\label{cor:r75-stationary-process}
For every stationary finite-alphabet process that is not i.i.d., the
minimum pure-start edge-emitting hidden Markov realization has exactly
one more label than the minimum freely mixed realization, with the same
infinite-value convention.  In particular, the $q$ versus $q+1$ conclusion
of Theorem~\ref{thm:r74-depolarizing-obstruction} is an instance of a
general stationary principle, while its phase-dependent $q$ remains a
separate spectral conclusion.
\end{corollary}

Equality of these complete response functions also preserves the law
under every common randomized policy for commands, stopping and terminal
readout.  For each finite stopped history, multiply the equal response by
the same policy probabilities.  Summation gives all finite stopping
events, and decreasing continuation events give the nontermination
probability.  Consequently all nonnegative costs measurable from that
record have equal expectations, including infinite expectations.  These
are process-realization label counts.  If the external policy has its own
persistent memory, that memory is counted separately in a closed-loop
implementation; it is not supplied free by the theorem.

\section{Approximation by truncating rare conditional excursions}
\label{sec:r76-rare-approximation}

Exact memory can be infinite although one finite simulator approximates
every stopped experiment with a prescribed expected cost.  The damping
construction of Theorem~\ref{thm:r75-stopped-approximation} gives such an
upper bound for the entire angle circle.  For a fixed finite command
alphabet we now obtain a smaller bound by preserving a long conditional
excursion exactly and modifying only its overflow.  The general mechanism
depends on the growth of the continuation semigroup, and applies beyond
the qubit family.

The ingredients have classical antecedents.  Polyhedral positive
realizations are part of the theory of Heller and of Monr\`as and Winter
\cite{Heller65,MonrasWinter14,MonrasWinter16}.  Martingale calculations
for pattern waiting times are classical \cite{Li80}.  The use made here
is a uniform approximation of complete adaptive stopped laws, with
an explicit memory bound from a finite return polytope and the number of
distinct continuation states.  It is separate from the exact noisy
realization problem of \cite{Fanizza24}.

\subsection{A sharp cost bound for a long conditional excursion}

Consider sequential calls, each producing a report declared either
\emph{continuing} or \emph{returning}.  Conditional on the entire past,
the chosen command and any private randomization, suppose the probability
of a continuing report is at most $r$, where $0<r<1$.
Let $\eta_L$ be the first call completing $L$ consecutive continuing
reports, with $L\ge1$.  A policy can stop initially or after any call;
its number of calls is $\tau$.

\begin{lemma}[Sharp expected-cost excursion bound]
\label{lem:r76-excursion-budget}
For every such policy with $\E\tau\le B<\infty$,
\begin{equation}\label{eq:r76-excursion-budget}
 \Prb\{\eta_L\le\tau\}
 \le \min\left\{1,\frac{(1-r)B}{r^{-L}-1}\right\}.
\end{equation}
The coefficient is optimal under the stated assumptions.  For independent
continuing indicators of probability $r$, the supremum over randomized
stopping policies of expected cost at most $B$ is exactly the right side.
\end{lemma}
\begin{proof}
Before the first completed run, let $K_t\in\{0,\ldots,L-1\}$ be the
current run length.  It increases by one on a continuing report and
returns to zero otherwise.  Absorb it at $L$ after the first completed
run, and put
\[
 u(k)=\frac{r^{-k}-1}{1-r}\quad(0\le k\le L).
\]
For $k<L$ and any conditional continuation probability $p\le r$,
\[
 p u(k+1)+(1-p)u(0)-u(k)
 \le r u(k+1)-u(k)=1.
\]
Stop this calculation at $\sigma_H=\min\{\tau,\eta_L,H\}$.
Summing the conditional inequalities, including the indicator that the
next call is actually made, gives
\[
 u(L)\Prb\{\eta_L\le\tau,\ \eta_L\le H\}
 \le \E u(K_{\sigma_H})\le\E\sigma_H\le B.
\]
Let $H\to\infty$.  Increasing events give
\eqref{eq:r76-excursion-budget}; no assumption about the expected duration
of another process is used.

For sharpness, take independent Bernoulli indicators of probability $r$.
The first completed run has finite expectation: disjoint blocks of length
$L$ succeed with probability $r^L$, so its expectation is at most
$L/r^L$.  In this case the displayed drift equals one.  Stopping at
$\eta_L\wedge H$ and then using bounded convergence for $u(K)$ and
monotone convergence for the duration proves
$\E\eta_L=u(L)$.  With probability
$\min\{1,B/u(L)\}$, selected before making any call, run this experiment
until $\eta_L$; otherwise stop immediately.  Its expected cost is at most
$B$ and its success probability equals the right side of
\eqref{eq:r76-excursion-budget}.
\end{proof}

Sharpness in this lemma concerns the probability of the exceptional
run.  It will not imply that every simulator must incur the approximation
error estimated by that probability.

\subsection{A general finite return polytope construction}

Let $\mathcal H$ be finite dimensional and let $\mathcal C$ be a closed convex set
of density matrices invariant under every positive-probability normalized
branch below.  The full state space is allowed.  For each command in a nonempty
finite set $A$, let $\{\mathcal I_{a,y}:y\in Y_a\}$ be a quantum
instrument, with finite $Y_a$.  Select a subset $R_a\subseteq Y_a$ of
continuing reports.  Suppose that there are density matrices
$\sigma_1,\ldots,\sigma_k$, with
$P=\operatorname{conv}\{\sigma_1,\ldots,\sigma_k\}$, such that
\begin{equation}\label{eq:r76-return-hypotheses}
 \begin{split}
 &\rho_0\in P,\qquad
   \sum_{y\in R_a}\tr\mathcal I_{a,y}(\rho)\le r<1,\\
 &\frac{\mathcal I_{a,y}(\rho)}{\tr\mathcal I_{a,y}(\rho)}\in P
 \quad(y\notin R_a,\ \tr\mathcal I_{a,y}(\rho)>0)
 \end{split}
\end{equation}
for every command $a$ and every $\rho\in\mathcal C$, with the seeds in $\mathcal C$.
On the full state space, the continuation bound is equivalently the finite
family of operator inequalities
$\sum_{y\in R_a}\mathcal I_{a,y}^*(I)\preceq rI$.
The return condition is a specified geometric hypothesis; it is not
inferred merely from mixing of the unobserved channel.  A prescribed
preparation in $\mathcal C$ but outside $P$ can be included by adding it to the seed list.

Write $\mathcal E=\{(a,y):a\in A,\ y\in R_a\}$ and $b=|\mathcal E|$.
For a word $w$ over $\mathcal E$, let $\mathcal I_w$ be its chronological
composition.  For $L\ge1$ define the finite set
\begin{equation}\label{eq:r76-orbit-dictionary}
 \mathcal D_L=
 \left\{\frac{\mathcal I_w(\sigma_j)}{\tr\mathcal I_w(\sigma_j)}:
  1\le j\le k,\ |w|<L,\ \tr\mathcal I_w(\sigma_j)>0\right\},
 \qquad G(L)=|\mathcal D_L|.
\end{equation}
The empty word is included.  Repeated states may be identified.

\begin{theorem}[Finite-memory approximation from return geometry]
\label{thm:r76-return-approximation}
Assume \eqref{eq:r76-return-hypotheses} with $0<r<1$.
For each $L\ge1$, there is a single $G(L)$-label mixed-start classical
model, chosen independently of the policy and of $B$, such that
\begin{equation}\label{eq:r76-return-approximation}
 \sup_{\pi:\,\E_{\rm target}^{\pi}\tau\le B}
 \|\mathcal L_{\rm target}^{\pi}-\mathcal L_L^{\pi}\|_{\rm TV}
 \le \min\left\{1,\frac{(1-r)B}{r^{-L}-1}\right\}.
\end{equation}
Here policies can choose commands, stopping and arbitrary terminal POVMs
from the visible record and their own randomness.  The complete stopped
record includes all commands and reports, the stopping time, the terminal
outcome and nontermination.  The bound also holds in trace distance for
the classical--quantum terminal record, with nontermination represented
by its scalar mass.  There is no additional entangled quantum register.

A counted pure start needs at most $G(L)+1$ labels.  In general
\begin{equation}\label{eq:r76-semigroup-growth}
 G(L)\le k\sum_{j=0}^{L-1}b^j.
\end{equation}
If the $b\ge1$ continuation maps commute pairwise, then
\begin{equation}\label{eq:r76-commuting-growth}
 G(L)\le k\binom{L+b-1}{b}.
\end{equation}
More generally any independently proved upper bound for the number of
distinct normalized states in \eqref{eq:r76-orbit-dictionary} can be used
in place of these bounds.  If that number is $O(L^d)$, the construction
uses $O((\log(B/\delta))^d)$ labels for fixed $r$ and
$B/\delta\to\infty$.
\end{theorem}
\begin{proof}
Use the states of $\mathcal D_L$ as hidden labels.  At a label $\rho$,
the probability of report $y$ on command $a$ is
$p_{a,y}(\rho)=\tr\mathcal I_{a,y}(\rho)$.  Zero probabilities require
no successor choice.  For a returning report, choose any convex
decomposition of its normalized output in the seed states.  For a
continuing report whose normalized output belongs to $\mathcal D_L$,
use that output label deterministically.  In the remaining case, send
the state to any one fixed seed.  Multiply each successor distribution
by $p_{a,y}(\rho)$ to obtain the report transition row.  All rows are
nonnegative and the sums over the reports are stochastic.  Use a seed
decomposition of $\rho_0$ for the initial distribution.  At label $\rho$,
the terminal decoder for any POVM $(E_z)_z$ is $\tr(E_z\rho)$.

To verify the error, first use all finite continuation words, without
the length restriction in \eqref{eq:r76-orbit-dictionary}.  This countable
refinement reproduces the target exactly.  Every returning branch is
expanded in the seeds; every continuing branch uses its exact normalized
image.  On each individual hidden label the weighted state identity is
exact, so induction gives the target subnormalized quantum state for
every command--report history.  Initial mixing, terminal POVMs and common
randomized policies preserve these identities.  This is only a proof
device, not an uncounted finite implementation resource.

Couple this exact refinement to the finite model by identical initial
seed choices and by identical reports and successor choices while the
states agree.  The coupling first fails only when a continuing image
leaves $\mathcal D_L$.  After a return the exact label is a seed, so a
product of fewer than $L$ continuing branches is in $\mathcal D_L$ by
definition.  The same is true initially.  Thus a failure requires a run
of at least $L$ consecutive continuing reports.  Identifying duplicate
states does not change this implication: their next normalized image
depends only on the state itself, and every shorter word image is still
present.

Before a failure the two visible histories, quantum states and policy
choices can be identical.  After a failure use arbitrary valid tails.
For a fixed finite horizon, the stopped classical--quantum states have
a common positive contribution from all paths with no failure.  If the
remaining mass is $e$, their trace distance is at most $e$, by positivity
and the triangle inequality.  The same assertion applies to the complete
record after any terminal POVM.  The target's probability of a failure
before stopping is at most $\Prb_{\rm target}^{\pi}\{\eta_L\le\tau\}$,
which is bounded by Lemma~\ref{lem:r76-excursion-budget}.  Its conditional
continuation bound follows from \eqref{eq:r76-return-hypotheses}, also
after conditioning on the virtual seed choices.

For unbounded stopping, retain every stopped history up to horizon $H$
and combine the unfinished histories into one symbol.  The sum of trace
norms on the finite stopped blocks increases to its full sum; the two
unfinished masses decrease to their respective nontermination masses.
Their absolute difference therefore converges to the nontermination
term.  The uniform finite-horizon bound passes to the limit, giving
\eqref{eq:r76-return-approximation}.  The budget is always evaluated under
the target.  The argument does not require the approximate model to have
expected duration at most $B$.

For a pure start append one initializer.  On the first actual command it
uses the initial mixture times the corresponding report matrix and
enters the core.  Its zero-call terminal decoder is that of $\rho_0$.
It is never entered again.  This proves the upper $G(L)+1$ without an
uncounted preliminary randomization.  If a separate finite policy uses
$s$ labels, the same construction gives at most $sG(L)+1$ labels for
the closed product implementation with a prescribed initial label.

There are $b^j$ words of length $j$, proving
\eqref{eq:r76-semigroup-growth}.  For commuting maps the product depends
only on a multi-index $m\in\mathbb N^b$ with $|m|<L$; there are
$\binom{L+b-1}{b}$ such multi-indices.  Normalization and deletion of zero
images cannot increase this number.  Finally it suffices to take
\[
 L\ge
 \max\left\{1,
 \left\lceil\frac{\log(1+(1-r)B/\delta)}{\log(1/r)}\right\rceil\right\}
 \quad(B>0,\ \delta>0).
\]
The asserted consequence of polynomial orbit growth follows.
\end{proof}

The construction preserves every finite history containing no run of
$L$ continuing reports, including its unnormalized terminal state.
It does not claim that its approximate response is stationary or that
the extra pure-start label is necessary.  When all continuation maps
have finite normalized orbit on the seeds, the full finite orbit instead
gives an exact model.  These are upper constructions, not a converse
classification of general positive realizations.

\subsection{A polylogarithmic simulator for finitely many commanded angles}

Apply the theorem to \eqref{eq:r75-noisy-instrument} with $\alpha=1$.
Let the finite known angle alphabet be
$\Theta=\{\theta_1,\ldots,\theta_a\}$, where $a\ge1$.
The continuation report is $0$, of probability $1/8$.
In the real Bloch plane put
\[
 v_0=(1,0),\qquad
 v_1=(-1/2,\sqrt3/2),\qquad
 v_2=(-1/2,-\sqrt3/2).
\]
Let $\sigma_j$ be the corresponding pure states.  Their triangle contains
the disk of radius $1/2$, and $I/2$ is their equal mixture.
Report $1$ returns to $\sigma_0=\tau$.  By
Lemma~\ref{lem:r75-full-instrument-bounds}, every normalized report-$2$
image has Bloch norm at most $5/11<1/2$, so it returns to this same
triangle, at every angle and from every real-plane state.
The real-plane density matrices form an invariant closed convex set:
the unitaries are real, and the other two branches preserve real symmetric
matrices.  Thus they supply the set $\mathcal C$ in the general theorem.

\begin{corollary}[Uniform finite-alphabet approximation]
\label{cor:r76-qubit-polylog}
For the undamped instrument at any fixed finite angle alphabet $\Theta$
and initial state $I/2$, there is, for every integer $L\ge1$, one classical
model with at most
\begin{equation}\label{eq:r76-qubit-count}
 3\binom{L+a-1}{a}
 \quad\hbox{mixed-start labels, or}\quad
 1+3\binom{L+a-1}{a}
 \quad\hbox{counted pure-start labels},
\end{equation}
whose complete stopped-law error satisfies
\begin{equation}\label{eq:r76-qubit-error}
 \sup_{\pi:\,\E_1^\pi\tau\le B}
 \|\mathcal L_1^\pi-\mathcal L_L^\pi\|_{\rm TV}
 \le \min\left\{1,\frac{7B}{8(8^L-1)}\right\}.
\end{equation}
The stronger stopped classical--quantum distance has the same bound.
The statements include adaptive choices among all supplied angles and
arbitrary terminal qubit POVMs.

For $B>0$ and $\delta>0$ it is enough to choose
\begin{equation}\label{eq:r76-qubit-depth}
 L=\max\left\{1,
   \left\lceil\log_8\left(1+\frac{7B}{8\delta}\right)\right\rceil\right\}.
\end{equation}
For a single commanded angle this uses at most $3L+1$ pure-start labels,
uniformly in that angle.  For each fixed finite alphabet it uses
$O_a((\log(B/\delta))^a)$ labels as $\delta/B\downarrow0$.
When the sine and cosine of each commanded angle are algebraic, every
state and transition coefficient of the displayed construction is
algebraic and is obtained by finitely many exact field operations.
\end{corollary}
\begin{proof}
The continuation maps are scalar multiples of commuting rotations, so
their normalized orbit dictionary consists of
\begin{equation}\label{eq:r76-qubit-dictionary}
 r_{j,m}=R_{m_1\theta_1+\cdots+m_a\theta_a}v_j,
 \qquad 0\le j<3,\quad m\in\mathbb N^a,\quad |m|<L.
\end{equation}
We may keep all these indices, even if some represent the same state.
Initialize uniformly on $(0,0),(1,0),(2,0)$.  From a label with Bloch
vector $r$, and command $\theta_i$, set
\[
 h=\tfrac18+\tfrac1{16}r_1,\qquad
 p_2=\tfrac78-h,\qquad
 u_2=-\tfrac18 R_{\theta_i}r-hv_0.
\]
Report $1$ has a row of mass $h$ into seed $(0,0)$.  Report $2$ has rows
into the three seeds of masses
\begin{equation}\label{eq:r76-explicit-return-rows}
 t_{2,j}=\frac{p_2+2v_j\mathbin{\cdot}u_2}{3},
 \qquad j=0,1,2.
\end{equation}
Indeed, $\sum_jv_j=0$ and
$\sum_jv_jv_j^{\mathsf T}=\tfrac32I_2$, so these masses sum to $p_2$
and their weighted Bloch vector is $u_2$.  Moreover,
$|u_2|/p_2\le5/11$ and $p_2\ge11/16$ give
\[
 t_{2,j}\ge\frac{p_2}{3}(1-10/11)\ge\frac1{48}>0.
\]
Thus the rows give the exact report-$2$ subnormalized state, with an
explicit positive margin.  On report $0$, whose mass is $1/8$, move
from $(j,m)$ to $(j,m+e_i)$ if $|m|<L-1$.  At $|m|=L-1$ send that mass
to the three seeds uniformly.  All report probabilities are unchanged
on each represented input state, and their sum is one.  Only the state
following an overflow zero report is modified.

The resulting table uses words with the same multi-index identified at
each depth and the same common-prefix proof as the theorem.  It keeps
the depth indices and may truncate at depth $L$ even when a periodic
phase makes that next physical state coincide with an earlier one.
The first possible modified state still occurs on the $L$th consecutive
zero.  Its state count
is \eqref{eq:r76-qubit-count}, and substitution of $r=1/8$ in
\eqref{eq:r76-return-approximation} gives
\eqref{eq:r76-qubit-error}.  Formula~\eqref{eq:r76-qubit-depth} makes
this bound at most $\delta$.  The label count and all error bounds are
uniform in the numerical values of the finite supplied angles, although
the transition table itself uses those values.  The explicit rotations
and \eqref{eq:r76-explicit-return-rows} prove the algebraic assertion.
\end{proof}

If $B=0$, a single initial decoder suffices, with no call and no error.
The damping upper in Theorem~\ref{thm:r75-stopped-approximation} remains
available and can be smaller for particular $B,\delta,a$; one may choose
the smaller of the two proved upper counts.  Its simultaneous
all-angle statement also remains valid.  In contrast, the finite dictionary
in \eqref{eq:r76-qubit-dictionary} is for a \emph{fixed finite} angle
alphabet; the bound is not asserted for a continuum of possible commands.

For every non-root phase, exact finite classical realization is still
impossible by Theorem~\ref{thm:r74-depolarizing-obstruction}.  At that same
fixed physical instrument the present bound is logarithmic in $B/\delta$
for a single angle.  Thus an inverse-square-root exact noise law is not a
lower bound on approximation of bounded-cost stopped experiments.
Section~\ref{sec:r76-physical-approximation} proves a matching logarithmic
law within the explicitly restricted class of physical dictionaries;
a matching lower for arbitrary classical terminal decoder tables remains
open.  The stopped-law bound also does not control the approximate
model's expected duration: unbounded costs
require a further tail or uniform-integrability hypothesis.  External
policy storage, table bits, arithmetic precision and exact sampling remain
separate resources.

\subsection{A matching logarithmic law for physical dictionaries}
\label{sec:r76-physical-approximation}

The preceding uppers use genuine density matrices as hidden dictionary
states.  Within that representation class a matching lower, including
the leading logarithmic constant, is possible.  This restriction is
substantive: an arbitrary classical hidden-state model can have terminal
decoder tables that are not Born probabilities of fixed density matrices.
The conclusion below does not impose that restriction on such a model
without proof.

Fix one commanded angle $\theta=2\pi\omega$ in the undamped instrument.
A \emph{physical-dictionary simulator} has finitely many labels with
fixed qubit density matrices $\sigma_i$, arbitrary stochastic report
transition tables, and terminal decoder
$d_E(i)=\tr(E\sigma_i)$ for every requested effect $E$.
The causal transition tables are fixed independently of any later
terminal-measurement request.
The transitions themselves need not implement the quantum instrument on
each dictionary state.  Denote by $N_{\rm phys}^{\rm m}(\delta;B)$ and
$N_{\rm phys}^{\rm p}(\delta;B)$ the least mixed- and counted-pure-start
label counts that approximate, in total variation to error at most
$\delta$, every common visible stopping and terminal-POVM policy whose
target expected calls are at most $B$.

\begin{theorem}[Sharp leading order for a fixed badly approximable phase]
\label{thm:r76-physical-logarithm}
Suppose $\omega$ is irrational and has a supplied constant $c>0$ such that
\begin{equation}\label{eq:r76-phase-separation}
 \|q\omega\|_{\mathbb R/\mathbb Z}\ge c/q
 \qquad(q\in\mathbb Z_{\ge1}).
\end{equation}
Necessarily $c\le1/2$.  Fix $B>0$ and put $\zeta=\min\{1,7B/8\}$.
For every integer $J\ge2$, if
\begin{equation}\label{eq:r76-physical-lower-condition}
 0<\delta\le\frac{\zeta}{8^J(1+J^2/c^2)},
\end{equation}
then
\begin{equation}\label{eq:r76-physical-lower}
 N_{\rm phys}^{\rm m}(\delta;B)\ge J,
 \qquad N_{\rm phys}^{\rm p}(\delta;B)\ge J+1.
\end{equation}

There is also a finite integer $K_\omega$, independent of $B,\delta$,
such that, with $L$ from \eqref{eq:r76-qubit-depth},
\begin{equation}\label{eq:r76-physical-upper}
 N_{\rm phys}^{\rm m}(\delta;B)\le K_\omega+L-1,
 \qquad N_{\rm phys}^{\rm p}(\delta;B)\le K_\omega+L.
\end{equation}
Consequently, for either start convention and every fixed $B>0$,
\begin{equation}\label{eq:r76-physical-log-limit}
 \lim_{\delta\downarrow0}
 \frac{N_{\rm phys}(\delta;B)}{\log_8(1/\delta)}=1.
\end{equation}
For $\omega=(\sqrt5-1)/2$, one may take $c=1/3$ and $K_\omega=5$.
All approximation lower bounds in this theorem are for physical
dictionaries with the stated Born terminal decoders.
\end{theorem}
\begin{proof}
For $n\ge0$, consider the report word $1\,0^n$ and the pure target state
$\rho_n$ it prepares.  Its Bloch vector is
$v_n=(\cos(2\pi n\omega),\sin(2\pi n\omega),0)$, with the coordinates
ordered along the real Bloch plane.  The word has probability
$8^{-(n+1)}$.  Use the following policy.  With probability $\zeta$ begin
the test; otherwise stop before any call.  During the test stop as soon
as a report differs from the specified prefix, or at its completion.
Apply the terminal POVM $(\rho_n,I-\rho_n)$ on successful completion.
The target expected number of calls is
\[
 \zeta\sum_{j=0}^{n}8^{-j}\le\zeta\,8/7\le B.
\]
Thus this is one of the policies in the definition of the approximation
error, for every $n$, without a horizon-dependent loss of the budget.

In the simulator let $A_n=\sum_i a_{n,i}\sigma_i$, $a_{n,i}\ge0$, be
the subnormalized terminal state on the event that this test is begun
and its complete prefix is observed.  Put
$p_n=\zeta 8^{-(n+1)}$.  Two events in the full stopped record give
\begin{equation}\label{eq:r76-prefix-mass-fidelity}
 |\tr A_n-p_n|\le\delta,
 \qquad \tr[(I-\rho_n)A_n]\le\delta.
\end{equation}
The second event has zero target probability.  If $p_n>\delta$, these
inequalities imply that some dictionary state with positive weight obeys
\[
 \tr(\rho_n\sigma_i)
 \ge 1-\frac{\delta}{p_n-\delta}.
\]
This conclusion is derived from joint events.  No error bound for a
rare conditional distribution has been assumed.

For $0\le n<J$, condition \eqref{eq:r76-physical-lower-condition} gives
\[
 \frac{\delta}{p_n-\delta}\le c^2/J^2.
\]
If $r_i\in\mathbb R^3$, $|r_i|\le1$, is the Bloch vector of the
corresponding physical dictionary state, then
\[
 |r_i-v_n|^2
 \le 2-2r_i\mathbin{\cdot}v_n
 =4[1-\tr(\rho_n\sigma_i)]\le4c^2/J^2.
\]
On the other hand, for $0\le m<n<J$, the elementary inequality
$\sin(\pi x)\ge2x$ on $[0,1/2]$ and
\eqref{eq:r76-phase-separation} give
\[
 |v_n-v_m|=2|\sin(\pi(n-m)\omega)|
 \ge4\|(n-m)\omega\|_{\mathbb R/\mathbb Z}
 \ge\frac{4c}{n-m}>\frac{4c}{J}.
\]
One dictionary vector cannot be within distance $2c/J$ of two such
points.  At least $J$ distinct labels are therefore necessary.

For a pure start, the initial label has its own physical density matrix
$\sigma_o$.  Stopping before any call and choosing any binary terminal
POVM is allowed at zero cost.  The error condition therefore implies
$\tfrac12\|\sigma_o-I/2\|_1\le\delta$, or $|r_o|\le2\delta$.
Its fidelity with any pure state is at most $1/2+\delta$.
Under \eqref{eq:r76-physical-lower-condition}, $J\ge2$ and $c\le1/2$
give $\delta<1/4$ and $c^2/J^2\le1/16$.  Thus the initial label cannot
be one of the $J$ labels whose fidelity is at least $15/16$.
This proves both inequalities in \eqref{eq:r76-physical-lower}.

For the improved upper, choose $K=K_\omega$ such that
\[
 P_K=\operatorname{conv}\{v_0,\ldots,v_{K-1}\}
 \supseteq\{x\in\mathbb R^2:|x|\le5/11\}.
\]
Such a finite $K$ exists for every irrational $\omega$.  Here is a direct
argument sufficient for this use.  Pigeonholing finitely many fractional
multiples of $\omega$ gives an integer $q>0$ with
$0<d=\|q\omega\|<1/4$.  The points with indices
$0,q,2q,\ldots,\lfloor1/d\rfloor q$ have successive circular gaps at
most $d$, in one orientation or the other.  Taking a prefix containing
all these indices, its inscribed polygon has every central gap less
than $\pi/2$, and hence inradius greater than $1/\sqrt2>5/11$.
The origin is in this polygon.  It is a valid seed polytope: report $1$
returns to $v_0$, and every report-$2$ output is in its contained disk.

For one angle, all continuation images of these $K$ seeds through depth
$L-1$ belong to
$\{v_0,\ldots,v_{K+L-2}\}$, a set of $K+L-1$ states.  Returning
branches are expanded in the original $K$ seeds.  No overflow can occur
before $L$ successive zero reports, even when a return chooses the last
seed.  Theorem~\ref{thm:r76-return-approximation} therefore proves
\eqref{eq:r76-physical-upper}.  The finite barycentric coefficients can
be chosen from any triangulation of $P_K$.  All labels, including the
optional initializer representing $I/2$, are physical states.

For fixed $B$, the upper has
$L=\log_8(1/\delta)+O_B(1)$.  Conversely,
\eqref{eq:r76-physical-lower-condition} holds for an integer
\[
 J=\log_8(1/\delta)
       -2\log_8\log_8(1/\delta)+O_{B,c}(1)
\]
chosen sufficiently small, as follows by taking logarithms of that
inequality.  Dividing the two bounds by $\log_8(1/\delta)$ proves
\eqref{eq:r76-physical-log-limit}.

Finally put $\omega=(\sqrt5-1)/2$, and let $p$ be a nearest integer to
$q\omega$.  Its conjugate is $\omega'=-(\sqrt5+1)/2$.  The product
$(q\omega-p)(q\omega'-p)=p^2+pq-q^2$ is a nonzero integer.  Therefore
\[
 |q\omega-p|\ge\frac1{|q\omega'-p|}
 \ge\frac1{\sqrt5q+1/2}>\frac1{3q},
\]
proving \eqref{eq:r76-phase-separation} with $c=1/3$.
The first five fractional orbit angles, in circular order, are
\[
 0,\quad 2\omega-1,\quad4\omega-2,\quad\omega,\quad3\omega-1.
\]
Their successive gaps are $2\omega-1$ or $2-3\omega$.
Since $3/5<\omega<5/8$, their maximum is $2\omega-1<1/4$.
This proves the required polygon inclusion with $K_\omega=5$.
\end{proof}

The final limit fixes $B>0$; it is not a uniform two-parameter asymptotic
as $B$ also varies.  For arbitrary classical decoder tables, the general
finite-alphabet upper remains valid, while the present packing lower
need not apply.  The badly approximable phase condition has an explicit
quadratic-irrational instance above.  It is not asserted for the separate
algebraic rotation $(3+4i)/5$ studied earlier in the paper.

The improved upper does apply to that original algebraic rotation, with
an explicit constant.  Write $\theta=\arg((3+4i)/5)\in(0,\pi/2)$.
The inequalities $\pi/6<\theta<\pi/3$ put the first seven orbit points
$0,\theta,\ldots,6\theta$ in their circular order with every gap less
than $\pi$.  The six consecutive edges have distance
$\cos(\theta/2)=2/\sqrt5$ from the origin.  The closing edge has
distance $-\cos(3\theta)=117/125$.  Both exceed $5/11$, so $K=7$ is
valid.  This gives $L+6$ mixed labels or $L+7$ pure labels at the error
in \eqref{eq:r76-qubit-error}.  All those orbit coordinates are rational.
Triangulating their rational polygon gives rational initial and returning
barycentric rows, so the resulting finite transition tables can be chosen
rational.  This is an explicit upper for that algebraic instrument, not
a claim about its unresolved Diophantine lower certificate.

\section{The causal task-order transfer theorem}\label{sec:central}
Suppose that every terminating state-stopped competitor reaches at least $B_0\ge1$ live labels under $\theta_*$. Such a certificate must concern actual physical command compatibility, not printed action names. Define
\begin{equation}\label{eq:widths}
\begin{array}{c|cc}
 i&m_-(M)&m_+(M)\\\hline
 \cp&M&\lfloor M/D_0\rfloor\\
 \ev&M&\lfloor M/D_0\rfloor\\
 \st&M-B_0&\lfloor(M-B)/D_0\rfloor .
\end{array}
\end{equation}
The event upper also requires $M\ge B$. The trivial state-stopped lower $B_0=1$ is valid for a procedure initially waiting to make a call.

\begin{theorem}[Affine task-order transfer with a paid reference and common construction]\label{thm:transfer}
Assume all the prepared-kernel, full-report, task-order, effort, occupation, resource and common-construction hypotheses of Section~\ref{sec:data}. Fix $N,M$, put $b_-=N/\kappa_-$, $b_+=N/\kappa_+$, and suppose $b_-\ge1$ and $m_-(M)\ge1$. Set
\[
 t_N=(1-A_N)_+/u_-.
\]
If $t_N>0$, every one of the three risks satisfies
\begin{equation}\label{eq:central-lower}
 R_i^\Theta(N,M)\ge
 \bigl[l_-t_N\Q_{m_-(M)}(b_-/t_N)-E_N\bigr]_+.
\end{equation}
If $t_N=0$, the asserted lower is zero, with no evaluation of $\Q$ at infinity. An empty legal class still has risk $+\infty$.

Suppose also $m_+(M)\ge1$, $b_+\ge1$, and, for event stopping, $M\ge B$. If
\begin{equation}\label{eq:beta}
 \beta_N=\frac{l_+b_+}{1+b_+a_+}\ge1,
\end{equation}
one common Borel program with actual expected cost at most $N$ in every model attains
\begin{equation}\label{eq:central-upper}
 R_i^\Theta(N,M)\le
 \frac{u_+(1+b_+a_+)}{l_+}\Q_{m_+(M)}(\beta_N)+b_+e_+.
\end{equation}
It commits almost surely. If $\beta_N<1$, the always-accept first-attempt construction instead gives $R_i^\Theta(N,M)\le u_+\Q_{m_+(M)}(1)+e_+$ at cost at most $\kappa_+\le N$.

The lower is against every legal competing strategy and uses one actual witness. The upper is one parameter-independent strategy for all models. The two effort constants need not agree. Neither direction assumes that the comparison experiment belongs to $\Theta$ or that a normalized predictive update is uniformly contracting.
\end{theorem}

\begin{corollary}[Scale-uniform relative transfer]\label{cor:relative}
Under Theorem~\ref{thm:transfer}, suppose $A_N=E_N=a_+=e_+=0$. If $l_+b_+\ge1$, then
\begin{equation}\label{eq:relative}
 \frac{l_-}{u_-}\Q_{m_-(M)}(u_-b_-)
 \le R_i^\Theta(N,M)
 \le\frac{u_+}{l_+}\Q_{m_+(M)}(l_+b_+).
\end{equation}
The factors $l_\pm,u_\pm$ may be fixed nonunit constants, independent of budget and resolution. In particular no term proportional to $b_\pm$ times a fixed absolute calibration error is required in this relative-order class.
\end{corollary}
\begin{proof}
Substitute $t_N=1/u_-$ and $\beta_N=l_+b_+$ in the theorem.
\end{proof}

\begin{corollary}[Separate exact checkpoint and autonomous laws]\label{cor:exact}
Suppose the certificates are exact: $l_\pm=u_\pm=1$, all defects vanish, $\kappa_-=\kappa_+=\kappa$, and exports are the identity. Then, for $N\ge\kappa$,
\[
 R_\cp^\Theta(N,m)=\Q_m(N/\kappa).
\]
The constructed autonomous upper uses $B+m$ state-stopped labels, or $\max(B,m)$ event-stopped labels. If the mandatory live count is sharp, $B_0=B$, then
\begin{align}
 R_\st^\Theta(N,B+m)&=\Q_m(N/\kappa)\quad(m\ge1),\label{eq:typed-exact}\\
 R_\ev^\Theta(N,M)&=\Q_M(N/\kappa)\quad(M\ge B).\label{eq:event-exact}
\end{align}
The subtraction by $B$ belongs only to the state-stopped model. Nonidentity exports require the expanded dictionary in \eqref{eq:widths}, unless an optimal dictionary is invariant under all exports. These exact stopping distinctions are inherited from the complete companion and are reproved here to fix the resource semantics.
\end{corollary}
\begin{proof}
The same lower and upper profile coincide in each asserted interface. Identity exports give $D_0=1$. The packing and converse widths are proved in Lemma~\ref{lem:roles}. No lower on live states is used to forbid event-stopped reuse.
\end{proof}

The familiar absolute-defect theorem is also a specialization: take $l_\pm=u_\pm=1$ while retaining $A_N,E_N,a_+,e_+$. Thus relative and absolute calibration are composed in one statement. The chosen metric is the common scored task throughout; it is not changed between the converse and the construction.

\section{Proof of the central theorem}
\begin{lemma}[Profile attainment]\label{lem:attainment}
For every finite $m\ge1$ and $b\ge1$, the minimum in \eqref{eq:profile} is attained. The profile is nonincreasing in $m,b$ and left-continuous in $b>1$.
\end{lemma}
\begin{proof}
Write a dictionary as an ordered tuple in the compact set $H^m$, permitting repetitions. The feasible densities are a weak-star closed subset of the radius-$b$ ball of $L^\infty(\nu)$: positivity, the upper bound and the mass equality are each closed under pairing with $L^1(\nu)$. They are therefore weak-star compact. As the tuple varies, its squared-distance function varies uniformly on $H$, since all points belong to a bounded set. Along a convergent product net, split the objective difference into this uniform variation and the weak-star pairing against the fixed bounded measurable loss at the limit tuple. This proves continuity of the objective and attainment. Nets avoid an unwarranted sequential compactness assumption. Monotonicity follows from inclusion of feasible sets. For left continuity take an optimizer $h$ at $b>1$ and let $h_\lambda=(1-\lambda)h+\lambda$. It has unit mass and cap $b_\lambda=(1-\lambda)b+\lambda<b$. Its objective tends to the optimum as $\lambda\downarrow0$, while monotonicity bounds the optimum at $b_\lambda$ between the optimum at $b$ and this objective.
\end{proof}

\begin{lemma}[Actual-history comparison measure]\label{lem:comparison}
Fix any robustly admissible competitor of cost at most $N$, with its fixed terminal dictionary $C$. There are finite nonnegative measures $\xi_c$ on $X$, $c\in C$, such that $\xi=\sum_c\xi_c$ obeys
\begin{align}
 \xi&\le (\E_*A_\tau)\nu\le b_-\nu,\label{eq:measure-cap}\\
 1&\le u_-\xi(X)+A_N,\label{eq:measure-mass}\\
 \E_*\|Z-\widehat r\|^2
 &\ge l_-\sum_{c\in C}\int d_c\,d\xi_c-E_N.
 \label{eq:measure-risk}
\end{align}
The measure $\xi$ is not asserted to be the terminal law of any comparison execution.
\end{lemma}
\begin{proof}
Adjoin an independent uniform random variable to the current report on both sides. This is current-call randomness, not a seed retained between calls; it is sufficient to represent the competitor's randomized output and stopping decision. A finite-output kernel has such a representation by cumulative probabilities. At a fixed reached productive history $h$ and chosen command $a$, the post-report decision then gives an acceptance indicator $q_h(y)$ and an output $r_h(y)\in C$ on acceptance. Give $r_h$ an arbitrary fixed value on rejection. The notation $y$ now includes the current independent random variable. Integrating the original certificates over that variable preserves their constants. Dependence on the entire earlier history is permitted in these functions.

For each Borel $D\subset X$, set
\begin{equation}\label{eq:xi-definition}
 \xi_c(D)=\E_*\sum_{j\le\tau:\ j\ {
m productive}}
 \int\one_{\{\pi_{H_{j-1},a_j}(y)\in D\}}
 q_{H_{j-1}}(y)\one_{\{r_{H_{j-1}}(y)=c\}}
 P^0_{H_{j-1},a_j}(dy,H).
\end{equation}
The outer expectation is the actual pre-call occupation at $\theta_*$. All terms are jointly measurable and nonnegative. Tonelli defines the measures, whose total is finite because $\E_*A_\tau\le b_-$. Summing over $c$ and dropping the acceptance indicator gives \eqref{eq:measure-cap}: each comparison row has mark law $\nu$.

The analogous sum with the actual rows instead of $P^0$ has mass exactly one, because the actual procedure commits exactly once almost surely. Apply \eqref{eq:lower-mass} at each actual reached history, then sum. The occupation certificate yields \eqref{eq:measure-mass}. Apply \eqref{eq:lower-task} in the same way to the disjoint actual commitment events. Their expected losses sum to the actual risk. The comparison contribution is $d_{r_h(y)}(\pi_{h,a}y)$ by \eqref{eq:full-moments}; this identity conditions on the full current report and therefore does not overlook extra information seen by the competitor. The accumulated loss defect is at most $E_N$, giving \eqref{eq:measure-risk}. One may first truncate the number of calls. Monotone convergence for the nonnegative contributions and the finite defect bounds justify the limit.

No actual past is replaced by a comparison past. In particular the competitor need not terminate, or obey a cost bound, if hypothetically run under the comparison experiment. This distinction is why the argument applies when that experiment is not an operating model.
\end{proof}

\begin{remark}[Why the full report is required]
A constant mark can conceal a visible fair bit which also equals the future audit. A comparison using only that mark has $f=0,w=1$, but a two-value decoder reading the bit predicts it with zero loss. The hypotheses therefore condition on the full report seen by the competitor; matching a compressed mark marginal alone cannot establish the lower.
\end{remark}

\begin{lemma}[Live roles, output roles and physical compatibility]\label{lem:roles}
A program with $\ell$ live roles and $t$ output roles has a state-stopped implementation with $\ell+t$ labels and an event-stopped implementation with $\max(\ell,t)$ labels, preserving every path cost and complete stopped scored law. Conversely an event-stopped $M$-label program has at most $M$ decoder values. A state-stopped $M$-label program with at least $B_0$ reached live labels has at most $M-B_0$ terminal decoder values.

For a physical live-state certificate, consider finitely many mandatory pre-call contexts $I$ with no free context input. Suppose the commands legal at a history of context $i$ are exactly $\mathcal U_i$, and every terminating competitor visits every context with positive probability under the same witness. This positive probability may depend on the competitor; no uniform lower bound over competitors is required. Let $b_{\rm phys}$ be the least size of a cover of $I$ by subsets $J$ with $\bigcap_{i\in J}\mathcal U_i\ne\varnothing$. Then at least $b_{\rm phys}$ live labels are required. This certificate is invariant under duplicate names and bijective stateless recodings of the same physical commands.
\end{lemma}
\begin{proof}
For state stopping take disjoint copies of the live and terminal role sets and type them zero and one. For event stopping embed the live and terminal role sets separately into one alphabet of size $\max(\ell,t)$. On continuing transitions use the live embedding; on commitment emit the final event and use the terminal embedding. A command row is read only on a continuing branch, while the decoder is read only after commitment. Sharing their label therefore causes no conflict and uses no continuing mode bit. Couple the original and recoded coins. Every finite stopped prefix and every path cost agree, and almost-sure termination gives equality of the stopped scored laws. A fixed decoder has one value per retained label. In state stopping the global halt map excludes each reached live label from the output set, proving its subtraction. Event stopping has no such exclusion.

A live label reached in contexts $J$ must have one command distribution supported on all the corresponding $\mathcal U_i$. Since $J$ is finite, the union of the exceptional null command sets is null; hence their intersection supports a probability and is nonempty. The sets of contexts served by the labels cover $I$, proving the bound. Conversely a nonempty intersection contains a command and hence a Dirac command distribution. Replacing names without changing executed kernels, legal domains, costs and reports does not change these intersections. A translator which itself remembers a hidden phase is not such a recoding; its memory must be charged.
\end{proof}

\begin{proof}[Proof of Theorem~\ref{thm:transfer}]
Fix a competing program. Lemma~\ref{lem:roles} bounds its decoder width by $m_-(M)$. Let $s=\xi(X)$ in Lemma~\ref{lem:comparison}. If $t_N>0$, then \eqref{eq:measure-mass} gives $s\ge t_N$. Scale the measure to $\xi'=(t_N/s)\xi$. It has mass $t_N$ and remains bounded above by $b_-\nu$. Thus $h=(d\xi'/d\nu)/t_N$ has unit mass and cap $b_-/t_N$. Nearest-dictionary projection and nonnegativity imply
\[
 \sum_{c\in C}\int d_c\,d\xi_c
 \ge\int [w+\dist(f,C)^2]\,d\xi
 \ge t_N\Q_{m_-(M)}(b_-/t_N).
\]
Together with \eqref{eq:measure-risk} this proves the asserted lower for this competitor. Its worst-model risk dominates its risk at $\theta_*$. Taking the infimum over competitors gives \eqref{eq:central-lower}, without exchanging an infimum and a supremum. Nonnegative loss gives the zero truncation and the stated case $t_N=0$.

For the upper assume \eqref{eq:beta}. Choose an attained optimizer $(C,h)$ at width $m_+(M)$ and cap $\beta_N$. Let $r_C(x)$ be the first nearest point to $f(x)$ in a fixed ordering of $C$. After each attempt of $\mathcal G$, accept $x$ with probability $h(x)/\beta_N$, export $T_Jr_C(x)$ and commit; otherwise return to the start routing label. All these operations are Borel. The choice $J$ is used within the current attempt; any part retained thereafter is charged through the live role or exported decoder, not a free seed. The exports are assumed jointly measurable and take at most $D_0|C|$ distinct values, so their finite decoder and transitions are Borel as well.

Conditional on every surviving attempt history, the acceptance probability is at least
\begin{equation}\label{eq:hazard}
 p=\frac{l_+}{\beta_N}-a_+=\frac1{b_+}>0.
\end{equation}
Consequently the probability of reaching attempt $j+1$ is at most $(1-p)^j$, by induction on conditional survival probabilities. In particular commitment is almost sure and the expected number of attempts is at most $1/p=b_+$. Each begun attempt has conditional expected cost at most $\kappa_+$, not merely a cost conditional on a favorable outcome. Summing the predictable attempt costs gives actual expected cost at most $\kappa_+b_+=N$.

The conditional unnormalized loss contributed by an attempt is at most $u_+\Q_{m_+(M)}(\beta_N)/\beta_N+e_+$. Sum its contributions on the disjoint terminal events, using the same survival bound. The result is at most
\[
 b_+\left[\frac{u_+}{\beta_N}\Q_{m_+(M)}(\beta_N)+e_+\right],
\]
which is \eqref{eq:central-upper}. The program has $B$ live roles and at most $D_0m_+(M)$ terminal roles. Lemma~\ref{lem:roles} gives exactly the sufficient widths in \eqref{eq:widths}. A checkpoint counts only the terminal width. No physical renewal law was used: the conditional hazard controls survival even if the hidden state after rejection differs arbitrarily from its former state.

If $\beta_N<1$, commit after the first attempt, and use an optimizer at cap one. The loss certificate with $q=1$ gives $u_+\Q_{m_+(M)}(1)+e_+$, and the whole-attempt cost is at most $\kappa_+$. This proves all remaining assertions.
\end{proof}

\section{Scope and further mathematical questions}\label{sec:r73-scope}\label{sec:r74-scope}\label{sec:r75-scope}\label{sec:r76-scope}
The finite-group convex-specification theorem is general over the stated
finite alphabets, finite groups, closed convex constraints and arbitrary
stochastic controllers. Its uniform approximation quantifier ranges over
all finite command words. The resulting finite optimization is not a
closed numerical formula for every group and target. The exact
small-contrast multiplicity law and normal-stabilizer convexity criterion
remain general within their stated classes. The complete mean--contrast
formula concerns the displayed symmetric-group family for every $n\ge5$;
it does not classify high contrast for all finite groups. Its matrix
synthesis theorem is an existence equivalence over parameter pairs,
not a positive-factorization theorem for every given feasible matrix.

The return-polytope theorem applies to finite-command instruments on an
explicit invariant physical state set, with a uniform conditional bound
on continuing branches and a finite return polytope. These hypotheses
are verified for the undamped qubit application. A fixed finite angle
alphabet gives the new polylogarithmic upper; the retained damping
construction works simultaneously for the entire angle circle. The
matching logarithmic lower and leading constant require fixed $B>0$,
a fixed badly approximable phase, and physical dictionaries with Born
terminal readouts. They do not apply to arbitrary classical decoder
tables without a further argument. The original algebraic phase has an
explicit rational-table logarithmic upper; its analogous Diophantine
lower certificate is not supplied here.

The stationary initialization theorem applies to arbitrary predictively
stationary finite-alphabet response laws and makes no quantum,
reversibility, finite-rank or irreducibility assumption. The exact noisy
instrument laws concern their stated family and spectral regimes.
One-dependence and stationarity are assertions about a fixed command;
adaptive command choices can transmit dependence. A private simulator
label is not an observation available to the original policy. The
approximating process need not satisfy the target's expected stopping
budget: total variation alone does not control an unbounded expectation.

The general-$M$ squared-sign minimax-risk frontier, unrestricted
approximate-decoder lower bounds, arbitrary nonrigid robust feedback and
tied boundary decisions, and the intermediate phase-type timing frontier
remain distinct mathematical problems. None is replaced by a solved
response-memory problem. Program length, arithmetic precision, working
storage and exact sampling costs also remain distinct from live labels.

The actual-history comparison theorem and the independent causal, reset,
discrimination and geometric results keep their original quantifiers.
Local geometric conclusions remain local. Independent A1, A2 and further
Theta construction objectives require their own proofs. The supplement
contains the complete technical arguments and independent developments;
all inherited mathematical bodies and the separate R58 companion have
traceable preserved locations. Historical drafts and source-recovery
records are available without becoming prerequisites for reading the
current theorem chain.

\renewcommand{\bibliofont}{\footnotesize\setlength{\baselineskip}{10.8pt}}
\begin{thebibliography}{99}
\bibitem{Blackwell} D. Blackwell, \emph{Equivalent comparisons of experiments}, Ann. Math. Statist. \textbf{24} (1953), 265--272. DOI: 10.1214/aoms/1177729032.
\bibitem{CGM97} J. A. Cuesta-Albertos, A. Gordaliza and C. Matr\'an, \emph{Trimmed $k$-means: an attempt to robustify quantizers}, Ann. Statist. \textbf{25} (1997), 553--576. DOI: 10.1214/aos/1031833664.
\bibitem{GG99} L. A. Garc\'ia-Escudero and A. Gordaliza, \emph{Robustness properties of $k$ means and trimmed $k$ means}, J. Amer. Statist. Assoc. \textbf{94} (1999), 956--969. DOI: 10.1080/01621459.1999.10474200.
\bibitem{SC} C. R. Shalizi and J. P. Crutchfield, \emph{Computational mechanics: pattern and prediction, structure and simplicity}, J. Statist. Phys. \textbf{104} (2001), 817--879. DOI: 10.1023/A:1010388907793.
\bibitem{BC} N. Barnett and J. P. Crutchfield, \emph{Computational mechanics of input--output processes: structured transformations and the $\varepsilon$-transducer}, J. Statist. Phys. \textbf{161} (2015), 404--451; arXiv:1412.2690.
\bibitem{AMTW} A. Ambainis, M. Mosca, A. Tapp and R. de Wolf, \emph{Private quantum channels}, Proc. 41st Annual Symposium on Foundations of Computer Science (2000), 547--553. DOI: 10.1109/SFCS.2000.892142.
\bibitem{Buscemi} F. Buscemi, \emph{Comparison of quantum statistical models: equivalent conditions for sufficiency}, Comm. Math. Phys. \textbf{310} (2012), 625--647. DOI: 10.1007/s00220-012-1421-3.
\bibitem{BR14} N. Bertschinger and J. Rauh, \emph{The Blackwell relation defines no lattice}, Proc. IEEE International Symposium on Information Theory (2014), 2479--2483. DOI: 10.1109/ISIT.2014.6875280; arXiv:1401.3146.
\bibitem{Banerjee05} A. Banerjee, X. Guo and H. Wang, \emph{On the optimality of conditional expectation as a Bregman predictor}, IEEE Trans. Inform. Theory \textbf{51} (2005), 2664--2669. DOI: 10.1109/TIT.2005.850145.
\bibitem{Stengel96} B. von Stengel, \emph{Efficient computation of behavior strategies}, Games Econom. Behav. \textbf{14} (1996), 220--246. DOI: 10.1006/game.1996.0050.
\bibitem{Balder89} E. J. Balder, \emph{On compactness of the space of policies in stochastic dynamic programming}, Stochastic Process. Appl. \textbf{32} (1989), 141--150. DOI: 10.1016/0304-4149(89)90058-6.
\bibitem{Barry14} J. Barry, D. T. Barry and S. Aaronson, \emph{Quantum partially observable Markov decision processes}, Phys. Rev. A \textbf{90} (2014), 032311. DOI: 10.1103/PhysRevA.90.032311; arXiv:1406.2858.
\bibitem{Ying18} S. Ying and M. Ying, \emph{Reachability analysis of quantum Markov decision processes}, Inform. Comput. \textbf{263} (2018), 31--51. DOI: 10.1016/j.ic.2018.09.001; arXiv:1406.6146.
\bibitem{Saldi25} N. Saldi, S. Sanjari and S. Y\"uksel, \emph{Quantum Markov decision processes: general theory, approximations, and classes of policies}, SIAM J. Control Optim. \textbf{63} (2025), 1156--1183. DOI: 10.1137/24M1643864; arXiv:2402.14649.
\bibitem{Seidenberg54} A. Seidenberg, \emph{A new decision method for elementary algebra}, Ann. of Math. (2) \textbf{60} (1954), 365--374. DOI: 10.2307/1969640.
\bibitem{Chatterjee15} K. Chatterjee, M. Chmel\'ik, R. Gupta and A. Kanodia, \emph{Optimal cost almost-sure reachability in POMDPs}, Proc. AAAI Conference on Artificial Intelligence \textbf{29} (2015). DOI: 10.1609/aaai.v29i1.9683.
\bibitem{Andriushchenko23} R. Andriushchenko, A. Bork, M. \v{C}e\v{s}ka, S. Junges, J.-P. Katoen and M. Mac\'ak, \emph{Search and explore: symbiotic policy synthesis in POMDPs}, Computer Aided Verification (2023); arXiv:2305.14149v2.
\bibitem{Halmos69} P. R. Halmos, \emph{Two subspaces}, Trans. Amer. Math. Soc. \textbf{144} (1969), 381--389. DOI: 10.1090/S0002-9947-1969-0251519-5.
\bibitem{Kessler21} D. A. Kessler, E. Barkai and K. Ziegler, \emph{First-detection time of a quantum state under random probing}, Phys. Rev. A \textbf{103} (2021), 022222. DOI: 10.1103/PhysRevA.103.022222; arXiv:2012.01763.
\bibitem{DelVecchio26} G. Del Vecchio Del Vecchio and S. N. Majumdar, \emph{Optimal detection of quantum states via projective measurements}, J. Phys. A: Math. Theor. \textbf{59} (2026), 035001. DOI: 10.1088/1751-8121/ae34fe; arXiv:2509.08556.
\bibitem{Companion} Q. Qi, \emph{General Theta Foundations I: Acquired Geometry and Causal Resource Transfer}, complete technical companion on absolute calibration, exact stopping resources and nonintervening sensing, 2026. The complete retained R58 native source and its independently built companion accompany this manuscript; their identities are fixed in the source-provenance and preservation records.
\bibitem{Tzeng92} W.-G. Tzeng, \emph{A polynomial-time algorithm for the equivalence of probabilistic automata}, SIAM J. Comput. \textbf{21} (1992), 216--227. DOI: 10.1137/0221017.
\bibitem{LiYuYing12} Y. Li, N. Yu and M. Ying, \emph{Termination of nondeterministic quantum programs}, arXiv:1201.0891 (2012).
\bibitem{Fu24} J. Fu, H. Jiang, M. Xu, Y. Deng and Z.-B. Li, \emph{Algorithmic analysis of termination problems for nondeterministic quantum programs}, arXiv:2402.15827 (2024).

\bibitem{LiuVerdu18} J. Liu and S. Verd\'u, \emph{Rejection sampling and noncausal sampling under moment constraints}, Proc. IEEE International Symposium on Information Theory (2018), 1565--1569. DOI: 10.1109/ISIT.2018.8437857.
\bibitem{GocFlamich24} D. Goc and G. Flamich, \emph{On channel simulation with causal rejection samplers}, Proc. IEEE International Symposium on Information Theory (2024). DOI: 10.1109/ISIT57864.2024.10619339; examined text arXiv:2401.16579v2 (3 May 2024).
\bibitem{LiLiu26} Z. Li and T. Liu, \emph{Ranking experiments under sequential sampling}, arXiv:2608.19897v2 (28 August 2026). Version 1 appeared under the title \emph{The order of binary experiments under endogenous stopping}, with Z. Li as sole author; the version-2 statement is the one compared here.
\bibitem{Ying14pre} S. Ying and M. Ying, \emph{Reachability analysis of quantum Markov decision processes}, arXiv:1406.6146 (2014), Theorem 3.6; published journal version listed above as \cite{Ying18}.

\bibitem{KaraYuksel20} A. D. Kara and S. Y\"uksel, \emph{Robustness to incorrect system models in stochastic control}, SIAM J. Control Optim. \textbf{58} (2020), 1144--1182. DOI: 10.1137/18M1208058; examined text arXiv:1803.06046v7.
\bibitem{Demirci25} Y. E. Demirci, A. D. Kara and S. Y\"uksel, \emph{Sensitivity of filter kernels and robustness bounds to transition and measurement kernel perturbations in partially observable stochastic control}, arXiv:2508.10658v3 (2025).
\bibitem{Watrous18} J. Watrous, \emph{The Theory of Quantum Information}, Cambridge University Press, 2018, Chapter 3, Corollary 3.40.
\bibitem{KRN65} K. Kuratowski and C. Ryll-Nardzewski, \emph{A general theorem on selectors}, Bull. Acad. Polon. Sci. S\'er. Sci. Math. Astronom. Phys. \textbf{13} (1965), 397--403.

\bibitem{Charnes63} A. Charnes, W. W. Cooper and K. O. Kortanek, \emph{Duality in semi-infinite programs and some works of Haar and Carath\'eodory}, Management Sci. \textbf{9} (1963), 209--228. DOI: 10.1287/mnsc.9.2.209.
\bibitem{Shapiro09} A. Shapiro, \emph{Semi-infinite programming, duality, discretization and optimality conditions}, Optimization \textbf{58} (2009), 133--161. DOI: 10.1080/02331930902730070; author preprint, Optimization Online (2008).
\bibitem{Dufour25} F. Dufour and T. Prieto-Rumeau, \emph{Absorbing Markov decision processes and their occupation measures}, SIAM J. Control Optim. \textbf{63} (2025), 676--698. DOI: 10.1137/24M1661959. Related author preprint examined: \emph{A note on weak compactness of occupation measures for an absorbing Markov decision process}, arXiv:2402.10672v2 (2024).
\bibitem{Franc23} V. Franc, D. Prusa and V. Voracek, \emph{Optimal strategies for reject option classifiers}, J. Mach. Learn. Res. \textbf{24} (2023), no. 11, 1--49.
\bibitem{Jencova16} A. Jen\v{c}ov\'a, \emph{Comparison of quantum channels and statistical experiments}, arXiv:1512.07016v2 (2016). The separately indexed shorter version arXiv:1601.06370 was withdrawn in favor of this extended text.
\bibitem{LiElGamal18} C. T. Li and A. El Gamal, \emph{Strong functional representation lemma and applications to coding theorems}, IEEE Trans. Inform. Theory \textbf{64} (2018), 6967--6978; arXiv:1701.02827v4.
\bibitem{AldousShepp87} D. Aldous and L. Shepp, \emph{The least variable phase type distribution is Erlang}, Comm. Statist. Stochastic Models \textbf{3} (1987), 467--473. DOI: 10.1080/15326348708807067.
\bibitem{Sion58} M. Sion, \emph{On general minimax theorems}, Pacific J. Math. \textbf{8} (1958), 171--176. DOI: 10.2140/pjm.1958.8.171.
\bibitem{PiunovskiyZhang24} A. Piunovskiy and Y. Zhang, \emph{Extreme occupation measures in Markov decision processes with an absorbing state}, SIAM J. Control Optim. \textbf{62} (2024), 65--90. DOI: 10.1137/23M1572398. Author text examined: \emph{Extreme occupation measures in Markov decision processes with a cemetery}, arXiv:2307.03158v1.
\bibitem{DufourMixtures} F. Dufour and T. Prieto-Rumeau, \emph{Absorbing Markov decision processes: geometric properties and sufficiency of finite mixtures of deterministic policies}, arXiv:2512.17511; examined HTML at version-1 URL, Theorems 4.3--4.4 and Proposition 4.5.
\bibitem{Mereghetti24} C. Mereghetti, B. Palano and P. Raucci, \emph{Unary quantum finite state automata with control language}, Appl. Sci. \textbf{14} (2024), no. 4, 1490. DOI: 10.3390/app14041490.
\bibitem{RabinScott59} M. O. Rabin and D. Scott, \emph{Finite automata and their decision problems}, IBM J. Res. Develop. \textbf{3} (1959), 114--125. DOI: 10.1147/rd.32.0114.
\bibitem{Jaeger00} H. Jaeger, \emph{Observable operator models for discrete stochastic time series}, Neural Comput. \textbf{12} (2000), 1371--1398. DOI: 10.1162/089976600300015411.
\bibitem{ThonJaeger15} M. Thon and H. Jaeger, \emph{Links between multiplicity automata, observable operator models and predictive state representations---a unified learning framework}, J. Mach. Learn. Res. \textbf{16} (2015), 103--147.
\bibitem{MonrasWinter14} A. Monras and A. Winter, \emph{Quantum learning of classical stochastic processes: the completely-positive realization problem}, LIPIcs \textbf{27}, TQC 2014 (2014), 99--109. DOI: 10.4230/LIPIcs.TQC.2014.99.

\bibitem{Heller65} A. Heller, \emph{On stochastic processes derived from Markov chains}, Ann. Math. Statist. \textbf{36} (1965), 1286--1291. DOI: 10.1214/aoms/1177700000.
\bibitem{CarlylePaz71} J. W. Carlyle and A. Paz, \emph{Realizations by stochastic finite automata}, J. Comput. System Sci. \textbf{5} (1971), 26--40. DOI: 10.1016/S0022-0000(71)80005-3.
\bibitem{Fanizza24} M. Fanizza, J. Lumbreras and A. Winter, \emph{Quantum theory in finite dimension cannot explain every general process with finite memory}, Comm. Math. Phys. \textbf{405} (2024), article 50. DOI: 10.1007/s00220-023-04913-4; arXiv:2209.11225. See Section 2.3 for the FRDN phase examples and Theorem 2 and Section 4.2 for the noise-dependent worst-phase classical-memory lower bound.
\bibitem{DmitrievDynkin46} N. Dmitriev and E. Dynkin, \emph{On characteristic roots of stochastic matrices}, Izv. Akad. Nauk SSSR Ser. Mat. \textbf{10} (1946), no. 2, 167--184. Math-Net.Ru: im3595 (Russian).
\bibitem{JohnsonPaparella17} C. R. Johnson and P. Paparella, \emph{A matricial view of the Karpelevi\v c theorem}, Linear Algebra Appl. \textbf{520} (2017), 1--15. DOI: 10.1016/j.laa.2017.01.009; arXiv:1611.06970v2.
\bibitem{KirklandSmigoc22} S. Kirkland and H. \v Smigoc, \emph{Stochastic matrices realising the boundary of the Karpelevi\v c region}, Linear Algebra Appl. \textbf{635} (2022), 116--138. DOI: 10.1016/j.laa.2021.11.016; arXiv:2110.01040v1.
\bibitem{BenvenutiFarina04} L. Benvenuti and L. Farina, \emph{A tutorial on the positive realization problem}, IEEE Trans. Automat. Control \textbf{49} (2004), 651--664. DOI: 10.1109/TAC.2004.826715.
\bibitem{BenvenutiFarina99} L. Benvenuti and L. Farina, \emph{An example of how positivity may force realizations of ``large'' dimension}, Systems Control Lett. \textbf{36} (1999), 261--266. DOI: 10.1016/S0167-6911(98)00098-X.
\bibitem{OCinneide90} C. A. O'Cinneide, \emph{Characterization of phase-type distributions}, Comm. Statist. Stochastic Models \textbf{6} (1990), 1--57. DOI: 10.1080/15326349908807134.
\bibitem{Maier91} R. S. Maier, \emph{The algebraic construction of phase-type distributions}, Comm. Statist. Stochastic Models \textbf{7} (1991), 573--602. DOI: 10.1080/15326349108807207.
\bibitem{HoneVarona19} A. N. W. Hone and J. L. Varona, \emph{Continued fractions and irrationality exponents for modified Engel and Pierce series}, Monatsh. Math. \textbf{190} (2019), 501--516. DOI: 10.1007/s00605-018-1244-1. Author version arXiv:1810.10127v1, equation (3.3).
\bibitem{Nussbaum98} R. D. Nussbaum and M. Scheutzow, \emph{Admissible arrays and a nonlinear generalization of Perron--Frobenius theory}, J. London Math. Soc. (2) \textbf{58} (1998), 526--544. DOI: 10.1112/S0024610798006796.
\bibitem{Brown64} D.~R. Brown, On clans of non-negative matrices,
\emph{Proc. Amer. Math. Soc.} \textbf{15} (1964), 671--674,
doi:10.1090/S0002-9939-1964-0165504-6.
\bibitem{Schwarz64} \v{S}. Schwarz, On the structure of the semigroup of
stochastic matrices, \emph{Publ. Math. Inst. Hungar. Acad. Sci.}
\textbf{9} (1964), 297--311, especially Theorems~7--8.
\bibitem{Flor69} P. Flor, On groups of non-negative matrices,
\emph{Compositio Math.} \textbf{21} (1969), 376--382,
\url{https://www.numdam.org/article/CM_1969__21_4_376_0.pdf}.
\bibitem{Serre07} J.-P. Serre, Bounds for the orders of the finite subgroups
of $G(k)$, in \emph{Group Representation Theory}, M. Geck, D. Testerman
and J. Th\'evenaz (eds.), EPFL Press, Lausanne, 2007, 405--450;
author lecture text, April 2006, archived as arXiv:1011.0346.
\bibitem{BabaiBealsRockmore93} L. Babai, R. Beals and D. Rockmore,
Deciding finiteness of matrix groups in deterministic polynomial time,
\emph{Proceedings of ISSAC 1993}, 117--126,
doi:10.1145/164081.164104.
\bibitem{KieferWachter14} S. Kiefer and B. Wachter, Stability and complexity
of minimising probabilistic automata, in \emph{ICALP 2014, Part II},
Lecture Notes in Computer Science \textbf{8573}, 268--279,
doi:10.1007/978-3-662-43951-7\_23; arXiv:1404.6673, Section~4.2.
\bibitem{Schwarz62} \v{S}. Schwarz, Homomorphisms of a completely simple
semigroup onto a group, \emph{Matematicko-fyzik\'alny \v{c}asopis}
\textbf{12} (1962), no.~4, 293--300;
\url{https://www.dml.cz/handle/10338.dmlcz/126333}.

\bibitem{GPT13} J. Gouveia, P. A. Parrilo and R. R. Thomas,
\emph{Lifts of convex sets and cone factorizations}, Math. Oper. Res.
\textbf{38} (2013), 248--264. DOI: 10.1287/moor.1120.0575;
arXiv:1111.3164, especially Theorem 2.12 on equivariant factorizations.
\bibitem{KPT10} V. Kaibel, K. Pashkovich and D. O. Theis,
\emph{Symmetry matters for sizes of extended formulations},
SIAM J. Discrete Math. \textbf{26} (2012), 1361--1382.
DOI: 10.1137/110839813; arXiv:0911.3712v3.
An extended abstract appeared in IPCO 2010, LNCS 6080, 135--148.
\bibitem{BocciCarliniRapallo11} C. Bocci, E. Carlini and F. Rapallo,
\emph{Perturbation of matrices and nonnegative rank with a view toward
statistical models}, SIAM J. Matrix Anal. Appl. \textbf{32} (2011),
1500--1512. DOI: 10.1137/110825455; arXiv:1010.3566.
\bibitem{GillisGlineur12} N. Gillis and F. Glineur,
\emph{On the geometric interpretation of the nonnegative rank},
Linear Algebra Appl. \textbf{437} (2012), 2685--2712.
DOI: 10.1016/j.laa.2012.06.038; arXiv:1009.0880.
\bibitem{Okay26} C. Okay, \emph{Polyhedral classical simulators for quantum
computation}, Res. Math. Sci. \textbf{13} (2026), article 55.
DOI: 10.1007/s40687-026-00626-x. In particular Theorem 16, Lemma 17
and Corollary 19 treat invariant-polytope updates and adaptive composition.
\bibitem{AubrunSzarek17} G. Aubrun and S. J. Szarek,
\emph{Dvoretzky's theorem and the complexity of entanglement detection},
Discrete Anal. (2017), paper no. 1, 20 pp.; arXiv:1510.00578,
especially Section 6 on polyhedral approximations of density matrices.
\bibitem{Fawzi21} H. Fawzi,
\emph{On polyhedral approximations of the positive semidefinite cone},
Math. Oper. Res. \textbf{46} (2021), 1479--1489.
DOI: 10.1287/moor.2020.1077; arXiv:1811.09649.
\bibitem{HolroydLiggett16} A. E. Holroyd and T. M. Liggett,
\emph{Finitely dependent coloring}, Forum Math. Pi \textbf{4} (2016),
e9, 43 pp. DOI: 10.1017/fmp.2016.7; arXiv:1403.2448.
See Section 5 and the finite hidden-Markov obstruction attributed there
to Schramm.
\bibitem{Serre77} J.-P. Serre, \emph{Linear Representations of Finite Groups},
Graduate Texts in Mathematics \textbf{42}, Springer-Verlag, New York, 1977.
DOI: 10.1007/978-1-4684-9458-7.

\bibitem{MonrasWinter16} A. Monras and A. Winter,
\emph{Quantum learning of classical stochastic processes: the
completely-positive realization problem}, J. Math. Phys. \textbf{57}
(2016), 015219. DOI: 10.1063/1.4936935; arXiv:1412.3634.

\bibitem{VuThanhGillisLecron23} O. Vu Thanh, N. Gillis and F. Lecron,
\emph{Bounded simplex-structured matrix factorization: algorithms,
identifiability and applications}, IEEE Trans. Signal Process.
\textbf{71} (2023), 2434--2447.
DOI: 10.1109/TSP.2023.3289704; arXiv:2209.12638v2,
especially Section IV-C and Theorem 3.
\bibitem{Gillis12} N. Gillis,
\emph{Sparse and unique nonnegative matrix factorization through data
preprocessing}, J. Mach. Learn. Res. \textbf{13} (2012), 3349--3386,
especially Section 4.3, Lemma 27, Example 2 and Theorem 29.
\url{https://jmlr.org/papers/v13/gillis12a.html}.

\bibitem{Rabin63} M. O. Rabin, \emph{Probabilistic automata},
Information and Control \textbf{6} (1963), 230--245.
DOI: 10.1016/S0019-9958(63)90290-0, especially Theorem 3 and Section VIII.

\bibitem{Li80}
S.-Y. R. Li, \emph{A martingale approach to the study of occurrence of
sequence patterns in repeated experiments}, Ann. Probab. \textbf{8}
(1980), no. 6, 1171--1176, doi:10.1214/aop/1176994578.

\end{thebibliography}

\end{document}
